
%%FILE blueprint/src/content.tex
% In this file you should put the actual content of the blueprint.
% It will be used both by the web and the print version.
% It should *not* include the \begin{document}
%
% If you want to split the blueprint content into several files then
% the current file can be a simple sequence of \input. Otherwise It
% can start with a \section or \chapter for instance.

This is a blueprint for a project which forms a part of the course
\emph{MS-EV0029 Introduction to Formalized Mathematics in Lean}
at Aalto University, running from February to June, 2025.
The goal is to formalize the Fisher-Tippett-Gnedenko theorem,
classifying univariate extreme value distributions.
Contributions are welcome by course participants!

\textbf{WARNING: The blueprint is still work in progress!}
Contributions by the course participants
to improve the blueprint and add details
(which are helpful for formalization) are very welcome!

\input{sections/cdf.tex}

\input{sections/extreme_value_distributions.tex}

\input{sections/classification.tex}

\input{sections/cdf_convergence.tex}

\input{sections/transforms.tex}

\input{sections/one_parameter_affine.tex}

\input{sections/types.tex}

\input{sections/lcinverse.tex}

\input{sections/cauchy_eqn.tex}


%%FILE blueprint/src/sections/cauchy_eqn.tex
\chapter{Cauchy-Hamel functional equation}

\section{Positive measure additive subgroups of the reals}

\begin{lemma}[Countably many connected components for an open set]
  \label{lem:countably-many-connected-components-of-open}
  %\uses{}
  \lean{IsOpen.countable_setOf_connectedComponentIn}
  \leanok
  Let $X$ be a locally connected separable space. Then any open subset $U \subseteq X$ has at
  most countably many connected components.
\end{lemma}
\begin{proof}
  % \uses{}
  \leanok
  (The proof is already formalized, see:
  \texttt{IsOpen.countable-setOf-connectedComponentIn}.)
\end{proof}

For subsets $A, B \subseteq \bR$, we use the following notation for pointwise
sum sets and difference sets:
\begin{align*}
  A + B & = \set{ a + b \; \big| \; a \in A, \, b \in B } , \\
  A - B & = \set{ a - b \; \big| \; a \in A, \, b \in B } .
\end{align*}

We denote by $\Leb$ the Lebesgue measure on $\bR$.

\begin{lemma}[Finding an interval with high overlap]
  \label{lem:exists-high-overlap-interval}
  %\uses{}
  \lean{exists_interval_measure_inter_gt_mul_measure}
  \leanok
  Let $A \subset \bR$ be a measurable set such that $0 < \Leb[A] < +\infty$.
  Then for any $r \in [0,1)$, there exists a nontrivial interval $J \subset \bR$
  (a subset of the real line which is connected and has nonempty interior)
  such that
  \begin{align*}
    \Leb[A \cap J] \; > \; r \, \Leb[J] .
  \end{align*}
\end{lemma}
\begin{proof}
  % \uses{}
  % \leanok
  Assume, without loss of generality, $0 < r < 1$.
  Since the Lebesgue measure is outer regular, we can find an open set $U \subset \bR$
  such that $A \subseteq U$ and $\Leb[U] < r^{-1} \, \Leb[A]$.

  The open set $U$ has at most countably many connected components
  (which are in fact open intervals); denote
  by $(U_i)_{i \in I}$ the indexed collection of them.

  Note that for at least one index $j \in I$ we
  have $\Leb[A \cap U_j] > r \Leb[U_j]$,
  because otherwise we get
  \begin{align*}
    \Leb[A] = \Leb[A \cap U]
    = \, & \sum_{i \in I} \Leb[A \cap U_i] \\
    \le \, & r \sum_{i \in I} \Leb[U_i] \\
    = \, & r \Leb[U],
  \end{align*}
  contradicting the choice of $U$.

  Now $J = U_j$ has the desired properties.
\end{proof}

\begin{lemma}[Shifts of a smaller interval contained in a larger interval]
  \label{lem:shift-interval-contained-in-larger}
  %\uses{}
  %\lean{}
  % \leanok % needs also `isPreconnected_of_Ioo_subset_of_subset_Icc` etc. before really done
  Let $I, J$ be nontrivial intervals, whose length satisfy
  \begin{align*}
    0 \; < \; \Leb[J] \; < \; \Leb[I] \; < \, +\infty .
  \end{align*}
  Then there exists an open interval $\Delta$ of length
  $\Leb[\Delta] = \Leb[I] - \Leb[J] > 0$ such that
  \begin{align*}
    t + J \; \subset \; I
    \qquad \text{ for any $t \in \Delta$.}
  \end{align*}
\end{lemma}
\begin{proof}
  % \uses{}
  % \leanok
  \ldots
\end{proof}

\begin{lemma}[Overlapping union of copies of an interval]
  \label{lem:overlapping-union-interval-copies}
  %\uses{}
  \lean{volume_union_add_self_ge_of_Ioo_subset, volume_union_add_self_le_of_subset_Icc,
    isConnected_of_Ioo_subset_of_subset_Icc}
  % \leanok % needs also `isPreconnected_of_Ioo_subset_of_subset_Icc` etc. before really done
  Let $J$ be a nontrivial interval of finite length ($0 \, < \, \Leb[J] \, < +\infty$).
  Let $c < 1$ and denote $\delta = c \Leb[J] \, < \, \Leb[J]$.
  Then for any $t \in (-\delta,\delta)$, the set
  $J' = (t+J) \cup J$ is an interval (connected set with nonempty interior)
  whose length satisfies the bound $\Leb[J'] \, < \, (1+c) \, \Leb[J]$.
\end{lemma}
\begin{proof}
  \uses{lem:shift-interval-contained-in-larger}
  % \leanok
  Denote $a = \inf J$ and $b = \sup J$. We have $-\infty < a < b < +\infty$ and
  \begin{align*}
    (a,b) \; \subseteq \; J \; \subseteq \; [a,b]
  \end{align*}
  and we have $\Leb[J] = b - a$.

  Let $t \in (+\delta,\delta)$, with $\delta = c \, \Leb[J] = c(b-a)$.
  Then we have the containments
  \begin{align*}
    \big( \min \{a,a+t\} , \max \{b,b+t\} \big)
      \; \subseteq \; (t+J) \cup J
      \; \subseteq \; \big[ \min \{a,a+t\} , \max \{b,b+t\} \big] .
  \end{align*}
  It in particular follows that $J' := (t+J) \cup J$ is an interval.

  If $t \ge 0$, then $J'$ is contained in
  $[a,b+t]$ which has length $b+t-a = (b-a) + t < \Leb[J] + c \, \Leb[J] = (1+c) \Leb[J]$.
  If $t < 0$, then $J'$ is contained in $[a+t,b]$ and one similarly gets a length bound.
\end{proof}


\begin{lemma}[Difference set of positive measure set contains an interval]
  \label{lem:difference-set-contains-interval}
  %\uses{}
  \lean{exists_Ioo_subset_diff_self_of_measure_pos}
  \leanok
  Let $A \subset \bR$ be a measurable set of positive Lebesgue measure.
  Then there exists a $\delta > 0$ such that
  \begin{align*}
    (-\delta,\delta) \, \subseteq \, A - A .
  \end{align*}
\end{lemma}
\begin{proof}
  \uses{lem:countably-many-connected-components-of-open, lem:exists-high-overlap-interval,
    lem:overlapping-union-interval-copies}
  % \leanok
  Pick a measurable subset $A_0 \subseteq A$ such that
  $0 < \Leb[A_0] < +\infty$.

  By Lemma~\ref{lem:exists-high-overlap-interval}, we can find a nontrivial interval~$J$
  such that $\Leb[A_0 \cap J] > \frac{3}{4} \Leb[J]$.
  Let $\delta = \frac{1}{2} \Leb[J]$.
  We claim that $(-\delta,\delta) \subseteq A_0 - A_0$ (which then
  clearly implies the assertion of the lemma).
  Indeed, suppose $t \in (-\delta,\delta)$.
  Then by Lemma~\ref{lem:overlapping-union-interval-copies},
  $(t+J) \cup J$ is an interval of length less than $\frac{3}{2} \Leb[J]$.
  Moreover, we have $A_0 \cap J \subseteq (t+J) \cup J$ and
  $t+(A_0 \cap J) \subseteq (t+J) \cup J$. Now note that
  \begin{align*}
    \Leb[t + (A_0 \cap J)] = \Leb[A_0 \cap J] > \frac{3}{4} \Leb[J] .
  \end{align*}
  If the sets $t + (A_0 \cap J)$ and $(A_0 \cap J)$ were disjoint, then
  the measure of $(A_0 \cap J) \cup (t+(A_0 \cap J))$ would thus be
  greater than $\frac{3}{2} \Leb[J]$, which is impossible given the length
  of the interval $(t+J) \cup J$ is less than $\frac{3}{2} \Leb[J]$.
  Therefore there exists a point $a \in (t + (A_0 \cap J)) \cap (A_0 \cap J)$.
  Denoting $a' = a - t$, we have $a, a' \in A_0 \cap J \subseteq A_0$
  and $t = a - a' \in A_0 - A_0 \subseteq A - A$. Since $t \in (-\delta,\delta)$
  was arbitrary, this proves the assertion.
\end{proof}

\begin{lemma}[Dividing a high overlap interval]
  \label{lem:dividing-high-operlap-interval}
  %\uses{}
  % \lean{}
  % \leanok
  Let $A$ be a measurable. Suppose that for some $a < b$,
  the interval $J = (a,b)$ satisfies
  $\Leb[A \cap J] \; > \;  r \, \Leb[J]$. %Denote $\delta = \Leb[J] = b-a$.
  Let $m \in \Npos$, and consider the subintervals
  \begin{align*}
    J_i = \Big( a + \frac{i}{(b-a)m} , \; a + \frac{i+1}{(b-a)m} \Big)
    \qquad \text{ for $i = 0, \ldots, m-1$.}
  \end{align*}
  Then for some $i$ we have $\Leb[A \cap J_i] \; > \;  r \, \Leb[J_i]$.
\end{lemma}
\begin{proof}
  % \uses{}
  % \leanok
\end{proof}

\begin{lemma}[Difference of two positive measure sets contains an interval]
  \label{lem:different-difference-set-contains-interval}
  % \uses{}
  % \lean{}
  % \leanok
  Let $A, B \subset \bR$ be two measurable sets of positive Lebesgue measure.
  Then the difference set $A - B$ contains a nontrivial open interval.
  %\begin{align*}
  %  (-\delta,\delta) \, \subseteq \, A - B .
  %\end{align*}
\end{lemma}
\begin{proof}
  \uses{lem:dividing-high-operlap-interval}
  % \leanok
  Pick measurable subsets $A_0 \subseteq A$ and $B_0 \subseteq B$ such that
  $0 < \Leb[A_0] < +\infty$ and $0 < \Leb[B_0] < +\infty$.

  By Lemma~\ref{lem:exists-high-overlap-interval}, we can find nontrivial intervals~$J, I$
  such that $\Leb[A_0 \cap J] > \frac{3}{4} \Leb[J]$
  and $\Leb[B_0 \cap I] > \frac{3}{4} \Leb[I]$.

  By Lemma~\ref{lem:dividing-high-operlap-interval}
  % we may divide
  % $I$ and $J$ to equal length subintervals, at least one of which inherits
  % this property, i.e.,
  for any $n,m \in \bN$ we can find intervals
  $I' \subset I$ and $J' \subset J$
  with $\Leb[I'] = \frac{1}{n} \Leb[I]$ and $\Leb[J'] = \frac{1}{m} \Leb[J]$
  and such that $\Leb[A_0 \cap J'] > \frac{3}{4} \Leb[J']$
  and $\Leb[B_0 \cap I'] > \frac{3}{4} \Leb[I']$.
  By choosing $n$ and $m$ suitably, we can ensure that
  \begin{align*}
    \frac{1}{2} \Leb[I'] \; \le \; \Leb[J'] \; < \; \Leb[I'] .
  \end{align*}

  Then by Lemma~\ref{lem:shift-interval-contained-in-larger}
  there exists an open interval $\Delta$ of length
  $\Leb[\Delta] \, > \, \Leb[I'] - \Leb[J']$ such that
  for all $t \in \Delta$ we have $t + J' \subset I'$.

  Consider a fixed $t \in \Delta$.
  Observe that $t + (B_0 \cap J) \subset I'$ and
  \begin{align*}
    \Leb[t + (B_0 \cap J')] \; = \; \Leb[B_0 \cap J']
      \; > \; \frac{3}{4} \Leb[J'] \; \ge \; \frac{3}{8} \Leb[I'] .
  \end{align*}
  We now claim that $A_0 \cap I'$ and $t + (B_0 \cap J')$ intersect.
  Their measures add up to at least
  \begin{align*}
    \Leb[A_0 \cap I'] + \Leb[t + (B_0 \cap J')]
    \; = \; & \Leb[A_0 \cap I'] + \Leb[B_0 \cap J'] \\
    \; > \; & \frac{3}{4} \Leb[I'] + \frac{3}{4} \Leb[J'] \\
    \; \ge \; & \frac{3}{4} \Leb[I'] + \frac{3}{8} \Leb[I']
    \; = \; & \frac{9}{8} \Leb[I'] .
  \end{align*}
  and both have been shown to be subsets of $I'$; therefore they
  cannot be disjoint.
  In particular there is a point $z \in (A_0 \cap I') \cap (t + B_0 \cap J')$,
  which means that
  \begin{align*}
    a = z = t + b
  \end{align*}
  for some $a \in A_0 \cap I'$ and $b \in B_0 \cap J'$. We solve
  $t = a - b \in A_0 - B_0 \subseteq A - B$. Since $t \in \Delta$
  was arbitrary, we have shown that the nontrivial open interval $\Delta$
  is contained in $A - B$.
\end{proof}


%%FILE blueprint/src/sections/cdf.tex
\chapter{Cumulative distribution functions}

\begin{definition}
  \label{def:cdf}
  \lean{CumulativeDistributionFunction}
  \leanok
  A function $F \colon \R \to \R$ is a
  \emph{cumulative distribution function} (c.d.f.) if
  \begin{enumerate}
  \item[(i)] $x \mapsto F(x)$ is increasing;
  \item[(ii)] $x \mapsto F(x)$ is right-continuous;
  \item[(iii)] $\lim_{x \to -\infty} F(x) = 0$ and $\lim_{x \to +\infty} F(x) = 1$.
  \end{enumerate}
\end{definition}

\begin{lemma}
  \label{lem:cdf-of-random-var}
  \uses{def:cdf}
  \lean{MeasureTheory.ProbabilityMeasure.cdf}
  \leanok
  If $X$ is a real-valued random variable, then the function
  $F \colon \R \to \R$ given by $F(x) = \PR \big[ X \le x \big]$ is a c.d.f.
\end{lemma}
\begin{proof}
  % \leanok
  Property (1.) in Definition~\ref{def:cdf} is obvious (by monotonicity of measures)
  and properties (2.) and (3.) are simple consequences of monotone convergence
  theorems for probability measures.
\end{proof}

\section{Degenerate distributions}

\begin{definition}
  \label{def:degenerate-cdf}
  \uses{def:cdf}
  \lean{CumulativeDistributionFunction.IsDegenerate}
  \leanok
  A c.d.f. $F$ is said to be
  \emph{degenerate} if for every $x \in \R$ we have either $F(x) = 0$ or $F(x) = 1$.
  Otherwise $F$ is said to be \emph{nondegenerate}.
\end{definition}

\begin{lemma}
  \label{lem:degenerate-cdf-iff-exists-jump}
  \uses{def:degenerate-cdf}
  \lean{CumulativeDistributionFunction.isDegenerate_iff}
  \leanok
  $F$ is a degenerate c.d.f. if and only if there exists a $x_0 \in \R$ such that
  \begin{align*}
  F(x) = \begin{cases}
          0 & \text{ for } x < x_0 \\
          1 & \text{ for } x \ge x_0 .
         \end{cases}
  \end{align*}
\end{lemma}
\begin{proof}
  \uses{}
  \leanok
  The ``if'' direction is clear. To prove the ``only if'' direction, assume that
  $F$ is a degenerated c.d.f., and let $x_0 = \inf \set{ x \in \R \; \big| \; F(x) = 1}$.
  Then it is straightforward to show by properties of a c.d.f. that $F$ has the asserted form.
\end{proof}

\begin{lemma}
  \label{lem:delta-has-degenerate-cdf}
  \uses{def:degenerate-cdf}
  \lean{CumulativeDistributionFunction.diracProba_is_degenerate}
  \leanok
  The c.d.f. of Dirac delta mass $\delta_{x_0}$ at $x_0 \in \R$ is degenerate.
\end{lemma}
\begin{proof}
  \uses{lem:degenerate-cdf-iff-exists-jump}
  \leanok
  \ldots
\end{proof}

\begin{lemma}
  \label{lem:degenerate-cdf-is-delta}
  \uses{def:degenerate-cdf}
  \lean{CumulativeDistributionFunction.eq_diracProba_of_isDegenerate}
  \leanok
  If a c.d.f. $F$ is degenerate, then it is the c.d.f. of a Dirac delta mass $\delta_{x_0}$
  at some point $x_0 \in \R$.
\end{lemma}
\begin{proof}
  \uses{lem:degenerate-cdf-iff-exists-jump}
  \leanok
  \ldots
\end{proof}

\section{Distributions of maxima of independent random variables}

% \begin{lemma}
% \label{lem:random-var-of-cdf}
% \uses{def:cdf}
% If $F \colon \R \to \R$ is a c.d.f., then there exists a probability
% space $(\Omega, \mathcal{F}, \PR)$ and a real-valued random variable
% $X \colon \Omega \to \R$ such that for every $x \in \R$
% we have $F(x) = \PR \big[ X \le x \big]$.
% \end{lemma}
% \begin{proof}
% Abstract nonsense (take $\Omega = [0,1]$ and $X = \idOf{\R}$ ).
% \end{proof}

\begin{lemma}
  \label{lem:cdf-of-max-two}
  \uses{def:cdf}
  % \leanok
  Let $X$ and $Y$ be two independent real-valued random variables with respective
  cumulative distribution functions $F$ and $G$,
  i.e. $F(x) = \PR \big[ X \le x \big]$ and $G(x) = \PR \big[ Y \le x \big]$.
  Then the c.d.f. of $M = \max (X, Y)$ is $x \mapsto F(x) \, G(x)$.
\end{lemma}
\begin{proof}
  \uses{}
  % \leanok
  Fix $x \in \R$. Note that $\max (X, Y) \le x$ if and only if both $X \le x$ and $Y \le x$.
  Calculate, using independence,
  \begin{align*}
  \PR\big[ \max (X, Y) \le x \big]
  \; = \; \PR\big[ X \le x \, \; Y \le x \big]
  \; = \; \PR\big[ X \le x \big] \; \PR\big[ Y \le x \big]
  \; = \; F(x) \, G(x) .
  \end{align*}
\end{proof}

\begin{lemma}
  \label{lem:cdf-of-max-many}
  \uses{def:cdf}
  % \leanok
  Let $X_0, X_1, \ldots, X_{n-1}$ be independent identically distributed real-valued random variables
  with cumulative distribution functions $F$,
  i.e. $F(x) = \PR \big[ X_j \le x \big]$ for every $j$.
  Then the c.d.f. of \[ M_n = \max_{0 \le j < n} X_j\] is the function $x \mapsto \big(F(x)\big)^n$.
\end{lemma}
\begin{proof}
  \uses{lem:cdf-of-max-two}
  % \leanok
  Induction on~$n$ using \ref{lem:cdf-of-max-two}.
\end{proof}

\begin{lemma}
  \label{lem:cdf-of-affine-max-many}
  \uses{def:cdf}
  % \leanok
  Let $X_0, X_1, \ldots, X_{n-1}$ be independent identically distributed real-valued random variables
  with cumulative distribution functions $F$,
  i.e. $F(x) = \PR \big[ X_j \le x \big]$ for every $j$,
  and let $a > 0$ and $b \in \R$.
  Then the c.d.f. of
  \begin{align*}
  \hat{M}_n = \frac{\max_{0 \le j < n} X_j - b}{a}
  \end{align*}
  is the function $x \mapsto \big(F(a x + b)\big)^n$.
\end{lemma}
\begin{proof}
  \uses{lem:cdf-of-max-many}
  % \leanok
  Use \ref{lem:cdf-of-max-many} and do a change of variables.
\end{proof}

\section{Distributions of minima of independent random variables}

\section{Equivalence classes modulo affine transformations}

\begin{definition}
  \label{def:oriented-affine-isomorphism}
  \lean{AffineIncrEquiv}
  \leanok
  The collection of all transformations $\R \to \R$ of the form
  $x \mapsto a x + b$, where $a>0$, $b \in \R$, forms a group.
  We call this the \emph{orientation preserving affine isomorphism group}
  and denote it by $\OriAffR$.
\end{definition}

\begin{definition}
  \label{def:oriented-affine-transform-of-cdf}
  \uses{def:oriented-affine-isomorphism, def:cdf}
  \lean{CumulativeDistributionFunction.affineTransform}
  \leanok
  The action of an orientation preserving affine isomorphism $A \in \OriAffR$
  on a cumulative distribution function $F$ is defined so
  that $A.F \colon \R \to \R$ is given by $(A.F)(x) = F(A^{-1}(x))$.
  Then $A.F$ is also a c.d.f.
\end{definition}

\begin{lemma}
  \label{lem:oriented-affine-action-on-cdf}
  \uses{def:oriented-affine-transform-of-cdf}
  \lean{CumulativeDistributionFunction.instMulActionAffineIncrEquiv}
  \leanok
  The actions of orientation preserving affine isomorphisms
  on a cumulative distribution functions is a group action, i.e.,
  $1.F = F$ and $(AB).F = A.(B.F)$ for any c.d.f. $F$ and any $A,B \in \OriAffR$.
\end{lemma}
\begin{proof}
  \uses{}
  % \leanok
  Direct calculations.
\end{proof}

\begin{lemma}
  \label{lem:degenerate-cdf-transform}
  \uses{lem:oriented-affine-action-on-cdf, def:degenerate-cdf}
  \lean{CumulativeDistributionFunction.affine_isDegenerate_iff}
  \leanok
  Let $F$ be a cumulative distribution function
  and $A \in \OriAffR$ an orientation preserving affine isomorphism.
  Then $A.F$ is degenerate if and only if $F$ is degenerate.
\end{lemma}
\begin{proof}
  \uses{}
  \leanok
  Straightforward from the definitions.
\end{proof}

\begin{lemma}
  \label{lem:cdf-continuity-pt-transform}
  \uses{lem:oriented-affine-action-on-cdf}
  \lean{CumulativeDistributionFunction.affine_continuousAt_of_continuousAt}
  \leanok
  Let $F$ be a cumulative distribution function,
  and $A \in \OriAffR$ an orientation preserving affine isomorphism.
  If a point $x \in \R$ is a continuity point of $F$, then
  the point $A(x) \in \R$ is a continuity point of $A.F$.
\end{lemma}
\begin{proof}
  \uses{}
  % \leanok
  Straightforward.
\end{proof}


% \begin{lemma}
% \label{lem:cdf-affine-group-action}
% \uses{def:cdf}
% Suppose $a > 0$ and $b \in \R$. Then for any c.d.f. $F$, also
% the function $x \mapsto F(a x + b)$ is a c.d.f.
% \end{lemma}
% \begin{proof}
% \ldots
% \end{proof}
%
% \begin{definition}
% \label{def:cdf}
% %\lean{CumulativeDistributionFunction}
% %\uses{def:degenerate-cdf}
% \uses{def:cdf}
% A function $F \colon \R \to \R$ is a
% \emph{cumulative distribution function} (c.d.f.) if
% \begin{enumerate}
% \item[(i)] $x \mapsto F(x)$ is increasing;
% \item[(ii)] $x \mapsto F(x)$ is right-continuous;
% \item[(iii)] $\lim_{x \to -\infty} F(x) = 0$ and $\lim_{x \to +\infty} F(x) = 1$.
% \end{enumerate}
% \end{definition}
%

\section{Miscellaneous results on cumulative distribution functions}

\begin{lemma}[Continuity points of c.d.f.s are those which carry no point mass]
  \label{lem:cdf-continuity-pt-iff-measure-singleton}
  \uses{def:cdf}
  \lean{CumulativeDistributionFunction.continuousAt_iff}
  \leanok
  Let $F$ be cumulative distribution function of a
  probability measure $\PRmeas$ on $\bR$. A point $x \in \bR$ is
  a continuity point of $F$ if and only if $\PRmeas[\set{x}] = 0$.
\end{lemma}
\begin{proof}
  % \uses{}
  % \leanok
  A c.d.f. is always continuous from the right.

  Continuity of $F$ from the left at $x$ means that for any
  sequence $(x_n)_{n \in \bN}$ increasing to $x$
  (i.e., $x_n \le x_{n+1} < x$ for all $n \in \bN$)
  \begin{align*}
    F(x_n) \to F(x) ,
  \end{align*}
  or equivalently in terms of measures
  \begin{align*}
    \PRmeas \big[ (-\infty,x_n] \big] \to \PRmeas \big[ (-\infty,x] \big] .
  \end{align*}
  But by monotone convergence of measures, we always have
  \begin{align*}
    \PRmeas \big[ (-\infty,x_n] \big]
    \to \; & \PRmeas \big[ (-\infty,x) \big] \\
    = \; & \PRmeas \big[ (-\infty,x] \big] - \PRmeas[\set{x}] .
  \end{align*}
  A comparison of these conditions shows that $F$ is also continuous from the left
  at $x$ if and only if $\PRmeas[\set{x}] = 0$.
\end{proof}

\begin{lemma}[A pair of nontrivial continuity points of nondegenerate c.d.f.]
  \label{lem:exists-two-nontrivial-continuity-pts-cdf}
  \uses{def:degenerate-cdf}
  \lean{CumulativeDistributionFunction.exists₂_continuousAt_of_not_isDegenerate}
  \leanok
  Let $G$ be a nondegenerate c.d.f. Then there exists
  continuity points $x_1 < x_2$ of $G$ such that
  $0 < G(x_1) \le G(x_2) < 1$.
\end{lemma}
\begin{proof}
  % \uses{}
  % \leanok
  Since $G$ is nondegenerate, there exists some $x_0 \in \bR$ such that
  $0 < G(x_0) < 1$. Since $G$ is continuous from the right, for some
  small $\delta > 0$ we have that $0 < G(x) < 1$ for
  all $x \in [x_0,x_0+\delta)$. Since the continuity points of $G$
  are dense, from any nonempty open interval we may pick a continuity point.
  First pick a continuity point $x_1 \in (x_0,x_0+\delta)$,
  and then pick another continuity point $x_2 \in (x_1,x_0+\delta)$.
\end{proof}


\begin{lemma}[Equality of c.d.f.s on a dense set suffices]
  \label{lem:cdf-equal-on-dense}
  \uses{def:cdf}
  \lean{CumulativeDistributionFunction.eq_of_forall_dense_eq}
  \leanok
  Suppose that $F,G$ are two c.d.f.s and $S \subseteq \bR$ is a dense
  subset of the real line. If $F(\xi) = G(\xi)$ for all $\xi \in S$,
  then we have $F = G$.
\end{lemma}
\begin{proof}
  \uses{}
  \leanok
  We must prove that for any $x \in \bR$ we have $F(x) = G(x)$.
  But by right-continuity of c.d.f.s, density of $S$, and coincidence of
  $F$ and $G$ on $S$, we have
  \begin{align*}
      F(x)
    = \; & \lim_{\xi \to x^+ \text{ along $S$}} F(\xi) \\
    = \; & \lim_{\xi \to x^+ \text{ along $S$}} G(\xi)
    = \, G(x) .
  \end{align*}
\end{proof}



\section{Topology on orientation-preserving affine isomorphisms}

\begin{definition}[]
  \label{def:affine-transform-topology}
  \uses{def:oriented-affine-isomorphism}
  \lean{instTopologicalSpaceAffineIncrEquiv}
  \leanok
  We equip the space $\OriAffR$ of orientation-preserving affine isomorphisms
  with the topology of pointwise convergence, i.e., with the
  coarsest topology which makes the evaluations $A \mapsto A(x)$ continuous
  for all $x \in \bR$.
\end{definition}

\begin{lemma}[The coefficients of affine map depend continuously on the map]
  \label{lem:affine-coefficients-continuous}
  \uses{def:oriented-affine-isomorphism}
  \lean{AffineIncrEquiv.continuous_coefs_fst, AffineIncrEquiv.continuous_coefs_snd}
  \leanok
  The coefficients $a$ and $b$ of an orientation-preserving affine isomorphism
  $A(x) = a x + b$ depend continuously on $A$.
\end{lemma}
\begin{proof}
  % \uses{}
  % \leanok
  We may first write $a = A(1) - A(0)$ and $b = A(0)$.
  These depend continuously on $A$,
  since the evaluations $A(1)$ and $A(0)$ do.
\end{proof}

\begin{lemma}[Metrizability of the topology on oriented affine isomorphisms]
  \label{lem:affine-metrizable}
  \uses{def:affine-transform-topology}
  \lean{AffineIncrEquiv.homeomorph_prod, AffineIncrEquiv.metrizableSpace}
  \leanok
  The topology of pointwise convergence makes $\OriAffR$
  homeomorphic to $\bR^2$, and in particular
  metrizable.
\end{lemma}
\begin{proof}
  \uses{lem:affine-coefficients-continuous}
  % \leanok
  The essential claim is that the function
  \begin{align*}
    \mathrm{cfs} \colon \OriAffR \to (0,+\infty) \times \bR
  \end{align*}
  obtained by mapping $A(x) = a x + b$ to its coefficients $(a,b)$ is a homeomorphism.
  (The homeomorphism to $\bR^2 = \bR \times \bR$ follows by combining with the
  homeomorphism $(a,b) \mapsto (\log a, b)$.)

  The continuity of $\mathrm{cfs}$ follows from Lemma~\ref{lem:affine-coefficients-continuous}.
  For the continuity of the inverse, we must only check that for any $x \in \bR$,
  its composition with the point evaluation at $x$ is continuous. But the composition
  is $(a,b) \mapsto a x + b$, and the continuity is clear.
\end{proof}

\begin{lemma}[Inversion of orientation preserving affine isomorphisms is continuous]
  \label{lem:affine-inversion-continuous}
  \uses{def:oriented-affine-isomorphism}
  \lean{AffineIncrEquiv.continuous_inv, AffineIncrEquiv.invHomeomorph}
  \leanok
  The map $A \mapsto A^{-1}$ is continuous on $\OriAffR$.
  In particular, inversion defines a homeomorphism of $\OriAffR$ to itself.
\end{lemma}
\begin{proof}
  \uses{lem:affine-coefficients-continuous}
  % \leanok
  Calculate and use Lemma~\ref{lem:affine-coefficients-continuous}.
\end{proof}

\begin{lemma}[The action of oriented affine transforms on c.d.f.s is continuous]
  \label{lem:action-on-cdf-continuous}
  \uses{lem:oriented-affine-action-on-cdf}
  \lean{CumulativeDistributionFunction.continuous_mulAction}
  \leanok
  The action $A . F$ of $A \in \OriAffR$ on a c.d.f $F$ depends
  jointly continuously on $A$ and $F$.

  \emph{(The topology on c.d.f.s is the topology of convergence in distribution, i.e.,
  convergence at all continuity points of the limit cdf.)}
\end{lemma}
\begin{proof}
  \uses{lem:affine-metrizable, lem:affine-inversion-continuous}
  % \leanok
  The spaces are metrizable, so it suffices to check sequential continuity.

  Suppose that $A_n \to B$ (oriented affine isomophisms)
  and $F_n \tendsinlaw G$ (c.d.f.s) as $n \to \infty$.
  Let $x \in \bR$ be a continuity point of $B.G$.
  Let $\eps > 0$. Note that $B^{-1}(x)$ is a continuity point of $G$.
  Then there exists a $\delta > 0$ such that $|G(y) - G\big(B^{-1}(x)\big)| < \frac{\eps}{2}$
  when $|y-B^{-1}(x)|<\delta$.

  By density of continuity points of $G$, pick continuity points
  $y_-,y_+$ such that
  \begin{align*}
    B^{-1}(x)-\delta \; < \; y_- \; < \; B^{-1}(x) \; < \; y_+ \; < \; B^{-1}(x)+\delta .
  \end{align*}
  Since $F_n \tendsinlaw G$, we have that
  $F_n(y_\pm) \to G(y_\pm)$ as $n \to \infty$,
  and in particular there exists some $N$ such that for $n \ge N$
  we have $\big| F_n(y_\pm) - G(y_\pm) \big| < \frac{\eps}{4}$ for both $y_\pm$.

  By Lemma~\ref{lem:affine-inversion-continuous}
  and $A_n \to B$, we get that $A_n^{-1} \to B^{-1}$, and in particular
  there exists some $N'$ such that for $n \ge N'$ we have
  \begin{align*}
    y_- < A_n^{-1}(x) < y_+ .
  \end{align*}

  For $n \ge \max \set{N,N'}$, we then have
  \begin{align*}
    (A_n.F_n)(x) \;
    = \; & F_n \big( A_n^{-1}(x) \big) \\
    \le \; & F_n (y_+) \\
    < \; & G(y_+) + \frac{\eps}{4} \\
    \le \; & G \big( B^{-1}(x) \big) + \frac{\eps}{2} + \frac{\eps}{4} \\
    = \; & (B.G) (x) + \frac{3\eps}{4} \\
    < \; & (B.G) (x) + \eps
  \end{align*}
  and similarly
  \begin{align*}
    (A_n.F_n)(x) \;
    = \; & F_n \big( A_n^{-1}(x) \big) \\
    \ge \; & F_n (y_-) \\
    > \; & G(y_-) - \frac{\eps}{4} \\
    \ge \; & G \big( B^{-1}(x) \big) - \frac{\eps}{2} - \frac{\eps}{4} \\
    = \; & (B.G) (x) - \frac{3\eps}{4} \\
    > \; & (B.G) (x) + \eps ,
  \end{align*}
  which together yield $\big| (A_n.F_n)(x) - (B.G)(x) \big| < \eps$.

  Since $x$ was an arbitrary continuity point of $B.G$ and
  $\eps > 0$ was arbitrary, this proves that $A_n.F \tendsinlaw B.G$.
\end{proof}

% lem:oriented-affine-action-on-cdf


%%FILE blueprint/src/sections/cdf_convergence.tex
\chapter{Convergence in distribution with cdfs}

This chapter provides a standard characterization of convergence in distribution
(weak convergence of probability measures) on the real line in terms of
cumulative distribution functions.

\section{Convergence in distribution}

Convergence in distribution for random variables can be defined when the random variables
take values in a topological space, and it amounts to the weak convergence of the
probability measures that are the laws of those random variables. In the special case
of real-valued random variables, or probability measures on the real line, the definition
reads:
\begin{definition}[Weak convergence of probability measures]
  \label{def:convergence-in-distribution}
  %\uses{}
  \lean{MeasureTheory.ProbabilityMeasure.instTopologicalSpace,
  MeasureTheory.ProbabilityMeasure.tendsto_iff_forall_lintegral_tendsto}
  \leanok
  \mathlibok
  A sequence $(\PRmeas_n)_{n \in \bN}$ of Borel probability measures on $\bR$ converges
  weakly to a Borel probability measure $\PRmeas$ on $\bR$ if for all bounded continuous
  functions $f \colon \bR \to [0,+\infty)$ we have
  \begin{align*}
    \lim_{n \to \infty} \int_{\bR} f(x) \, \ud \PRmeas_n(x)
    = \int_{\bR} f(x) \, \ud \PRmeas(x) .
  \end{align*}
\end{definition}

\section{Auxiliary results}

\begin{lemma}[Monotone real functions have only countably many points of discontinuity]
  \label{lem:monotone-discontinuities-countable}
  %\uses{}
  \lean{Set.Countable.dense_compl, Monotone.countable_not_continuousAt}
  \leanok
  \mathlibok
  A monotone function $f \colon \bR \to \bR$ can have at most countably
  many points of discontinuity. In particular the set $D \subset \bR$
  of continuity points of $f$ is dense in $\bR$.
\end{lemma}
\begin{proof}
  %\uses{}
  \leanok
  \mathlibok
  (The proof should already be in Mathlib.)
\end{proof}

\begin{lemma}[Tightness of a cumulative distribution function]
  \label{lem:cdf-tight}
  \uses{def:cdf}
  \lean{CumulativeDistributionFunction.forall_pos_exists_lt_gt_continuousAt}
  \leanok
  Let $F$ be a cumulative distribution function. Then for any
  $\eps > 0$ there exists points $a,b \in \bR$ with $a < b$ such that
  $F(b) - F(a) > 1 - \eps$ and $F$ is continuous at the points $a$ and $b$.
\end{lemma}
\begin{proof}
  \uses{def:cdf, lem:monotone-discontinuities-countable}
  % \leanok
  Cumulative distribution functions satisfy
  $F(x) \downarrow 0$ as $x \downarrow - \infty$ and
  $F(x) \uparrow 1$ as $x \uparrow + \infty$. The
  required large difference $F(b) - F(a)$ is
  obtained by choosing $a$ small enough so that $F(a) < \frac{\eps}{2}$
  and $b$ large enough so that $F(b) > 1 - \frac{\eps}{2}$.
  In order to guarantee that $a < b$ and
  that $a$ and $b$ are continuity points of $F$, we
  recall that continuity points of the monotone function $F$ are dense
  by Lemma~\ref{lem:monotone-discontinuities-countable}, so we
  may decrease $a$ and increase $b$ as appropriate.
\end{proof}

\begin{lemma}[Subdivision with small mesh and within dense set]
  \label{lem:subdivision-dense}
  %\uses{}
  \lean{forall_exists_subdivision_diff_lt_of_dense}
  \leanok
  Let $D \subset \bR$ be a dense set and $a,b \in D$ with $a < b$.
  Then for any $\delta > 0$ there exists a $k \in \bN$
  and $a = c_0, c_1, \ldots, c_{k-1}, c_k = b \in D$ such that
  $|c_j - c_{j-1}| < \delta$ for all $j=1,\ldots,k$.
\end{lemma}
\begin{proof}
  %\uses{}
  % \leanok
  \ldots
\end{proof}

\begin{lemma}[Subdivision for continuous function approximation]
  \label{lem:continuous-function-approximation-subdivision}
  %\uses{}
  \lean{forall_exists_subdivision_dist_apply_lt_of_dense_of_continuous}
  \leanok
  Let $D \subset \bR$ be a dense set, let $f \colon \bR \to \bR$ be continuous,
  let $a, b \in D$ with $a < b$, and let $\eps > 0$.
  Then there exists a $k \in \bN$ and points
  $a=c_0 < c_1 < \cdots < c_{k-1} < c_k = b$ such that for
  each $j = 1, \ldots, k$ we have $c_j \in D$ and
  \begin{align*}
  \big| f(x) - f(c_j) \big| < \eps
  \qquad \text{ for } \qquad x \in [c_{j-1} , c_j] .
  \end{align*}
\end{lemma}
\begin{proof}
  \uses{lem:subdivision-dense}
  % \leanok
  On the compact interval
  $[a,b] \subset \bR$, the continuous function~$f$ is uniformly continuous, so for some
  $\delta > 0 $ we have $|f(x)-f(y)|<\eps$ whenever $|x-y|<\delta$ and $x,y \in
  [a,b]$.
  Now apply Lemma~\ref{lem:subdivision-dense} to choose $k$ and points
  $a=c_0 < c_1 < \cdots < c_{k-1} < c_k$ such that $c_j - c_{j-1} < \delta$
  and $c_j \in D$ for all $j = 1 , \ldots , k$.
  Now for any $j = 1 , \ldots , k$, since for $x \in [c_{j-1} , c_j]$ we have
  $|x - c_j| < \delta$, we get
  \begin{align*}
  \big| f(x) - f(c_j) \big| < \eps
  \end{align*}
  as desired.
\end{proof}

\begin{lemma}[Simple function integral as linear combination of cdf differences]
  \label{lem:simple-integral-cdf-difference}
  \uses{def:cdf}
  \lean{CumulativeDistributionFunction.integral_sum_indicator_eq}
  \leanok
  Let $a = c_0 < c_1 < \cdots < c_k = b$ and consider the
  linear combination of indicator functions
  \begin{align*}
    h(x) = \sum_{j=1}^k \alpha_j \; \indof{(c_{j-1},c_j]}(x) .
  \end{align*}
  Then the integral of $h$ with respect to a Borel probability measure $\PRmeas$ on $\bR$ whose
  can be written as
  \begin{align*}
    \int_{\bR} h(x) \, \ud \PRmeas (x) = \sum_{j=1}^k \alpha_j \, \big( F(c_j) - F(c_{j-1}) \big) ,
  \end{align*}
  where $F$ is the c.d.f. of $\PRmeas$.
\end{lemma}
\begin{proof}
  % \uses{}
  % \leanok
  \begin{align*}
  \int_{\bR} h \, \ud \PRmeas
   = \; & \int_{\bR}
          \Big( \sum_{j=1}^k \alpha_j \, \indof{(c_{j-1},c_j]}(x) \Big) \, \ud \PRmeas(x) \\
   = \; & \sum_{j=1}^k \alpha_j \; \int_{\bR} \indof{(c_{j-1},c_j]}(x) \, \ud \PRmeas(x) \\
   = \; & \sum_{j=1}^k \alpha_j \; \PRmeas \big[ (c_{j-1},c_j] \big] \\
   = \; & \sum_{j=1}^k \alpha_j \, \big( F_n(c_j) - F_n(c_{j-1}) \big)
  \end{align*}
\end{proof}

\begin{lemma}[One of the portmanteau implications]
  \label{lem:portmanteau-convergence-for-borel}
  \uses{def:convergence-in-distribution}
  \lean{MeasureTheory.ProbabilityMeasure.tendsto_measure_of_null_frontier_of_tendsto',
  MeasureTheory.ProbabilityMeasure.tendsto_measure_of_null_frontier_of_tendsto}
  \leanok
  \mathlibok
  Weak convergence of probability measures implies that if the boundary of a Borel set carries
  no probability mass under the limit measure, then the limit of the measures of the set equals
  the measure of the set under the limit probability measure.

  In other words, if $\lim_{n \to \infty} \PRmeas_n = \PRmeas$ in the sense of weak convergence
  of measures, Definition~\ref{def:convergence-in-distribution}, and if $A \subset \bR$
  is a Borel set such that $\PRmeas[\partial A] = 0$, then
  \begin{align*}
    \lim_{n \to \infty} \PRmeas_n [ A ] = \PRmeas [ A ] .
  \end{align*}
\end{lemma}
\begin{proof}
  % \uses{}
  \leanok
  \mathlibok
  (The proof is in Mathlib.)
\end{proof}

\section{Convergence in distribution from pointwise convergence of cdfs}

\begin{theorem}[Sufficient condition for convergence in distribution with cdfs]
  \label{thm:convergence-in-distribution-with-cdf}
  \uses{def:convergence-in-distribution, def:cdf}
  \lean{tendsto_of_forall_continuousAt_tendsto_cdf}
  \leanok
  Let $F$ and $F_n$, $n \in \bN$, be cumulative distribution functions of
  probability measures $\PRmeas$ and $\PRmeas_n$, $n \in \bN$, respectively, i.e.,
  \begin{align*}
      F(x) = \; & \PRmeas \big[(-\infty,x]\big] & & \text{for $x \in \bR$} \\
      F_n(x) = \; & \PRmeas_n \big[(-\infty,x]\big] & & \text{for $x \in \bR$ and $n \in \bN$.}
  \end{align*}
  If $\lim_{n \to \infty} F_n(x) = F(x)$ for all continuity points $x$ of $F$,
  then $\lim_{n \to \infty} \PRmeas_n = \PRmeas$ in the sense of weak convergence
  of measures, Definition~\ref{def:convergence-in-distribution}.
\end{theorem}
\begin{proof}
  \uses{lem:monotone-discontinuities-countable, lem:cdf-tight,
  lem:continuous-function-approximation-subdivision,
  lem:simple-integral-cdf-difference}
  % \leanok
  Let $D \subset \bR$ denote the set of continuity points of $F$. By
  Lemma~\ref{lem:monotone-discontinuities-countable}, $D$ is dense in $\bR$.
  Assume that $\lim_{n \to \infty} F_n(x) = F(x)$ for all $x \in D$.

  Let~$\eps > 0$. Choose, by Lemma~\ref{lem:cdf-tight}, points
  $a,b \in D$, $a<b$, such that $F(b) - F(a) > 1 - \eps$.

  Observe also that since $\lim_{n \to \infty} F_n(a) = F(a)$
  and $\lim_{n \to \infty} F_n(b) = F(b)$, there exists some $N_1$ such that
  we have
  \begin{align*}
  F_n(b) - F_n(a) > 1 - 2\eps \qquad \text{ for all } n \geq N_1 .
  \end{align*}

  Let $f \colon \bR \to \bR$ be bounded and continuous.
  By Lemma~\ref{lem:continuous-function-approximation-subdivision}
  we can choose points
  $a=c_0 < c_1 < \cdots < c_{k-1} < c_k = b$ such that
  for all $j = 1 , \ldots , k$ we have $c_j \in D$ and
  \begin{align*}
  \big| f(x) - f(c_j) \big| < \eps
  \qquad \text{ for } \qquad x \in [c_{j-1} , c_j] .
  \end{align*}
  Define the simple function $h \colon \bR \to \bR$ by
  \begin{align*}
  h(x) = \sum_{j=1}^k f(c_j) \; \indof{(c_{j-1},c_j]}(x)
  \end{align*}
  The above estimate shows that $|f(x) - h(x)| < \eps$ for all $x \in [a,b]$.
  By boundedness of~$f$, there exists a constant $K>0$ such that $|f(x)| \leq K$
  for all $x \in \bR$. Since $h$ vanishes outside $(a,b]$,
  the triangle inequality for integral with respect to $\PRmeas_n$ gives
  \begin{align*}
  \Big| \int_{\bR} f \, \ud \PRmeas_n - \int_{\bR} h \, \ud \PRmeas_n \Big|
  \, \leq \; & \, \underbrace{\int_{(a,b]} |f-h| \, \ud \PRmeas_n}_{\leq \eps}
      + \underbrace{\int_{\bR \setminus (a,b]} |f| \, \ud \PRmeas_n}_{%
      \leq K \, \PRmeas_n\big[ \bR \setminus (a,b] \big]} .
  \end{align*}
  When $n \geq N_1$, we have
  $\PRmeas_n\big[ \bR \setminus (a,b] \big]
  = 1 - \PRmeas_n\big[ (a,b] \big] = 1 - (F_n(b) - F_n(a)) < 2 \eps$,
  and thus the triangle inequality implies
  \begin{align*}
  \Big| \int_{\bR} f \, \ud \PRmeas_n - \int_{\bR} h \, \ud \PRmeas_n \Big|
  \, \leq \; & \eps + K \, 2 \eps = (1 + 2K) \, \eps .
  \end{align*}
  Similarly, integrating now with respect to $\PRmeas$ instead, one shows that
  \begin{align*}
  \Big| \int_{\bR} f \, \ud \PRmeas - \int_{\bR} h \, \ud \PRmeas \Big|
  \, \leq \; & (1 + K) \, \eps .
  \end{align*}
  It remains to consider the integrals of the function~$h$ with respect to both
  $\PRmeas_n$ and $\PRmeas$. By Lemma~\ref{lem:simple-integral-cdf-difference},
  these integrals are expressible in terms of the cumulative distribution functions,
  \begin{align*}
  \int_{\bR} h \, \ud \PRmeas_n
  = \; & \sum_{j=1}^k f(c_j) \, \big( F_n(c_j) - F_n(c_{j-1}) \big)
  \end{align*}
  and
  \begin{align*}
  \int_{\bR} h \, \ud \PRmeas
  % = \; & \sum_{j=1}^k f(c_j) \; \PRmeas_n \big[ (c_{j-1},c_j] \big]
  %      \\
  = \; & \sum_{j=1}^k f(c_j) \, \big( F(c_j) - F(c_{j-1}) \big) .
  \end{align*}
  The difference of the integrals of~$h$ with respect to these two can
  therefore be estimated as
  \begin{align*}
  \Big| \int_{\bR} h \, \ud \PRmeas - \int_{\bR} h \, \ud \PRmeas_n \Big|
  \, = \; & \, \Big| \sum_{j=1}^k f(c_j) \,
      \big( F(c_j) - F_n(c_j) - F(c_{j-1}) +  F_n (c_{j-1}) \big) \Big| \\
  \leq \; & \, \sum_{j=1}^k |f(c_j)| \; \Big(
      \big| F(c_j) - F_n(c_j) \big| + \big| F(c_{j-1}) +  F_n (c_{j-1}) \big|
      \Big) \\
  \leq \; & \, 2 k K \max_{j = 0 , \ldots , k} \big| F(c_j) - F_n(c_j) \big| .
  \end{align*}
  By our assumption~\textup{(ii)}, we have
  $\lim_{n \to \infty} F_n(c_j) = F(c_j)$ for each $j = 1 , \ldots, k$, so
  there exists $N_2$ such that for $n \geq N_2$ we have
  $\max_{j = 1 , \ldots , k} | F(c_j) - F_n(c_j) | < \frac{\eps}{k}$, and
  thus
  \begin{align*}
  \Big| \int_{\bR} h \, \ud \PRmeas - \int_{\bR} h \, \ud \PRmeas_n \Big|
  \leq \; & \, 2 K \eps .
  \end{align*}
  Combining the estimates we have obtained, for $n \geq \max(N_1 , N_2)$, we have
  \begin{align*}
  & \Big| \int_{\bR} f \, \ud \PRmeas - \int_{\bR} f \, \ud \PRmeas_n \Big| \\
  \leq \; \, &
      \underbrace{
      \Big| \int_{\bR} f \, \ud \PRmeas - \int_{\bR} h \, \ud \PRmeas \Big|}_{
          \leq (1+K) \eps}
      + \underbrace{
      \Big| \int_{\bR} h \, \ud \PRmeas - \int_{\bR} h \, \ud \PRmeas_n \Big|}_{
          \leq 2 K \eps}
      + \underbrace{
      \Big| \int_{\bR} h \, \ud \PRmeas_n - \int_{\bR} f \, \ud \PRmeas_n	\Big|}_{
          \leq (1+2K) \eps } \\
  \leq \; \, & (2 + 5 K) \eps .
  \end{align*}
  Since $\eps > 0$ was arbitrary, this shows that
  $\int f \, \ud \PRmeas_n \to \int f \, \ud \PRmeas$ as $n \to \infty$, so we
  have established the weak convergence $\PRmeas_n \to \PRmeas$
  according to Definition~\ref{def:convergence-in-distribution}.
\end{proof}

\begin{lemma}[Necessary condition for convergence in distribution with cdfs]
  \label{lem:cdf-convergence-from-convergence-in-distribution}
  \uses{def:convergence-in-distribution, def:cdf}
  \lean{CumulativeDistributionFunction.tendsto_apply_of_tendsto_of_continuousAt}
  \leanok
  Let $\PRmeas$ and $\PRmeas_n$, $n \in \bN$, be Borel
  probability measures on $\bR$, and let $F$ and $F_n$, $n \in \bN$,
  be their cumulative distribution functions, respectively, i.e.,
  \begin{align*}
      F(x) = \; & \PRmeas \big[(-\infty,x]\big] & & \text{for $x \in \bR$} \\
      F_n(x) = \; & \PRmeas_n \big[(-\infty,x]\big] & & \text{for $x \in \bR$ and $n \in \bN$.}
  \end{align*}
  If $\lim_{n \to \infty} \PRmeas_n = \PRmeas$ in the sense of weak convergence
  of measures, Definition~\ref{def:convergence-in-distribution}, then
  for all continuity points $x$ of $F$ we have
  $\lim_{n \to \infty} F_n(x) = F(x)$.
\end{lemma}
\begin{proof}
  \uses{lem:cdf-continuity-pt-iff-measure-singleton, lem:portmanteau-convergence-for-borel}
  \leanok
  Let $x \in \bR$ be a continuity point of $F$. Then we have $\mu [\set{x}] = 0$,
  by Lemma~\ref{lem:cdf-continuity-pt-iff-measure-singleton}.
  Note that the boundary of the Borel set $(-\infty,x] \subset \bR$ is
  the singleton $\partial (-\infty,x] = \set{x}$.
  Therefore the assumption $\lim_{n \to \infty} \PRmeas_n = \PRmeas$ implies that
  \begin{align*}
    \PRmeas_n \big[ (-\infty,x] \big] \to \PRmeas \big[ (-\infty,x] \big] ,
  \end{align*}
  by a general fact (Lemma~\ref{lem:portmanteau-convergence-for-borel})
  about weakly converging sequences of measures
  that a for Borel sets whose boundary carries no mass in the limit measure.
  In terms of the c.d.f.s, the above reads
  \begin{align*}
    F_n(x) \to F(x)
  \end{align*}
  as asserted.
\end{proof}


%%FILE blueprint/src/sections/classification.tex
\chapter{Classification of extreme value distributions}

\section{Auxiliary classification I}

\begin{lemma}[Second order differential equation for Q]
  \label{lem:ode-of-order-two-for-Q}
  %\uses{}
  \lean{ExtremeValueProject.ode_of_order_two_for_Q}
  \leanok
  Suppose that
  \begin{align*}
  Q \colon \bR \to \bR
  \end{align*}
  is differentiable and satisfies
  \begin{align*}
  Q(0) = 0 \qquad \text{ and } \qquad Q'(0) = 1
  \end{align*}
  and
  \begin{align*}
  Q(h+s) = Q(h) \alpha(s) + Q(s)
  \end{align*}
  for some $\alpha \colon \bR \to \bR$ and every $s,h \in \bR$.
  Then $Q$ is twice continuously differentiable and satisfies
  \begin{align}
    Q''(s) = Q'(s) \, Q''(0) \qquad \text{ for every } s \in \bR .
  \end{align}
\end{lemma}
\begin{proof}
  %\uses{}
  % \leanok
  Note first that the equation implies (rearranging and dividing by $h$),
  for any $s$ and $h \ne 0$,
  \begin{align*}
    \frac{Q(h+s) - Q(s)}{h} = \frac{Q(h)}{h} \alpha(s) .
  \end{align*}
  Taking the limit as $h \to 0$ and using $Q(0)=0$ yields
  $Q'(s) = Q'(0) \, \alpha(s) = \alpha(s)$, where we also
  took into account $Q'(0) = 1$. Therefore necessarily
  $\alpha = Q'$, and the equation can be rewritten in the form
  \begin{align*}
  Q(h+s) = Q(h) \, Q'(s) + Q(s) \qquad \text{ for } s,h \in \bR .
  \end{align*}

  Rearranging the equation, we find
  \begin{align*}
    Q(h+s) - Q(s) = Q(h) \, Q'(s)
  \end{align*}
  and interchanging the role of $s$ and $h$ also
  \begin{align*}
    Q(h+s) - Q(h) = Q(s) \, Q'(h) .
  \end{align*}
  Subtracting the last two equations yields
  \begin{align*}
    Q(s) - Q(h) = Q(s) \, Q'(h) - Q(h) \, Q'(s) ,
  \end{align*}
  which by rearranging and dividing by $h \ne 0$ yields
  \begin{align*}
     \frac{Q(h)}{h} \big( Q'(s) - 1 \big) = Q(s) \frac{Q'(h) - 1}{h} .
  \end{align*}
  Taking the limit $h \to 0$, recalling $Q(0)=1$ and $Q'(0)=1$,
  gives the existence of the second derivative $Q''(0)$ and the equation
  \begin{align*}
    Q'(s) - 1 = Q(s) Q''(0) .
  \end{align*}
  Solving $Q'(s) = 1 + Q''(0) \, Q(s)$ and recalling that $Q$
  is differentiable shows that $Q'$ is also differentiable, so
  $Q$ is indeed twice differentiable. Differentiating, we get
  the asserted equation
  \begin{align*}
    Q''(s) = Q'(s) \, Q''(0) .
  \end{align*}
  Since $Q'$ is differentiable and in particular continuous, this also shows
  that $Q''$ is continuous, i.e., that $Q$ is twice continuously differentiable.
\end{proof}

\begin{lemma}[Solution for Q]
  \label{lem:solve-Q}
  %\uses{}
  \lean{ExtremeValueProject.solve_Q}
  \leanok
  Suppose that $Q \colon \bR \to \bR$ is twice continuously differentiable
  and $Q'$ is positive and $Q$
  and satisfies $Q(0)=0$, $Q'(0) = 1$, and the equation concluded in
  Lemma~\ref{lem:ode-of-order-two-for-Q} with $\gamma = Q''(0)$, i.e.,
  \begin{align}\label{eq: Q eqn dd}
  Q''(s) = \gamma \, Q'(s) \qquad \text{ for every } s \in \bR .
  \end{align}
  Then $Q$ is given by
  \begin{align*}
  Q(s) = \begin{cases}
    \frac{e^{\gamma s} - 1}{\gamma} & \; \text{ if $\gamma \ne 0$} \\
    \;\; s & \; \text{ if $\gamma = 0$}
    \end{cases}
    \qquad \text{for $s \in \bR$.}
  \end{align*}
%  if $\gamma \ne 0$ and
%  \begin{align*}
%  Q(s) = s
%  \end{align*}
%  if $\gamma = 0$.
  \end{lemma}
\begin{proof}
  %\uses{}
  % \leanok
  Since $Q$ is differentiable and $Q'(s)>0$ for any $s \in \bR$,
  we can write \eqref{eq: Q eqn dd} as
  \begin{align*}
  \der{s} \log Q'(s)
    = \frac{Q''(s)}{Q'(s)}
    = \gamma .
  \end{align*}
  Integrating and noting $\log Q'(0) = \log 1 = 0$, this
  yields $\log Q'(s) = \gamma s$ for all $s \in \bR$, or equivalently
  $Q'(s) = e^{\gamma s}$. Integrating a second time
  and noting $Q(0) = 0$, this yields the desired formulas
  in the two cases $\gamma \ne 0$ and $\gamma = 0$, respectively, as follows.

  If $\gamma \ne 0$, then we get
  \begin{align*}
  Q(s) = \int_0^s e^{\gamma u} \, \ud u = \frac{e^{\gamma s} - 1}{\gamma} .
  \end{align*}

  If $\gamma = 0$, then we get
  \begin{align*}
  Q(s) = \int_0^s e^0 \, \ud u = s .
  \end{align*}
\end{proof}

% TODO: Move to a different file.
\begin{theorem}[Monotone functions are a.e. differentiable]
  \label{thm:monotone-ae-differentiable}
  %\uses{}
  \lean{Monotone.ae_differentiableAt}
  \leanok
  \mathlibok
  If $f : \bR \to \bR$ is nondecreasing, then the derivative
  $f'(x)$ exists at almost every $x \in \bR$.
  In particular there exists points $x$ where $f'(x)$ exists.
\end{theorem}
\begin{proof}
  %\uses{}
  \leanok
  \mathlibok
  (The proof is already in Mathlib.)
\end{proof}

\begin{lemma}[Solution for E]
  \label{lem:solve-E}
  %\uses{lem:solve-Q}
  \lean{ExtremeValueProject.solve_E}
   \leanok
  Suppose that $E \colon (0,\infty) \to \bR$
  is nondecreasing and nonconstant function which satisfies $E(1) = 0$ and
  \begin{align*}
  E(\lambda \sigma) = E(\lambda) A(\sigma) + E(\sigma)
  \end{align*}
  for some $A \colon (0,\infty) \to (0,\infty)$ and all $\lambda, \sigma > 0$.
  Then, denoting $c = E'(1)$ and $\gamma = \frac{1}{E'(1)} \dder{s} E(e^s) \big|_{s=0}$,
  for all $\lambda \in \bR$ we have
  \begin{align*}
  E(\lambda) = \begin{cases}
    c \big( \lambda^\gamma - 1 \big) & \;\;\text{ if $\gamma \ne 0$.} \\
    c \, \log (\lambda) & \;\;\text{ if $\gamma = 0$.}
    \end{cases}
  \end{align*}
\end{lemma}
\begin{proof}
  \uses{lem:ode-of-order-two-for-Q, lem:solve-Q, thm:monotone-ae-differentiable}
  % \leanok
  Denote $H(s) = E(e^s)$ for $s \in \bR$. Then $H(0) = E(1) = 0$ and
  $H$ is also nondecreasing and nonconstant.
  By Theorem~\ref{thm:monotone-ae-differentiable}, there
  exists some $s_0 \in \bR$ such that the derivative $H'(s_0)$ exists.

  The equation for $E$ yields an equation for $H$,
  \begin{align*}
  H(h+s) = H(h) A(e^{s}) + H(s) \qquad \text{ for any } s, h \in \bR.
  \end{align*}
  We can rearrange this equation and divide by $h$,
  and we get for any $s, h \in \bR$, $h \ne 0$,
  \begin{align*}
  \frac{H(h+s) - H(s)}{h} = \frac{H(h)}{h} A(e^{s}) .
  \end{align*}
  If $s = s_0$, then the LHS tends to $H'(s_0)$ as $h \to 0$.
  The RHS must therefore also have a limit as $s \to 0$, and
  observing that $\frac{H(h)}{h} = \frac{H(h) - H(0)}{h}$,
  that limit is $H'(0) A(e^{s_0})$, which shows that the derivative $H'(0)$ exists.
  Then applying the same equation at general $s \in \bR$ shows that
  $H'(s) = H'(0) A(e^{s})$ must exist, so $H \colon \bR \to \bR$ is
  in fact everywhere differentiable.

  Note that we must have $H'(0) > 0$, because if $H'(0) = 0$
  then by the above equation $H'(s) = 0$ for every $s$ and then
  $H$ is a constant function, which is a contradiction.
  Denote $c = H'(0) = \der{s} E(e^s) \big|_{s=0} = E'(1)$.

  Also the equation and positivity of the function $A$ give $H'(s)>0$
  for all $s$. %, so $H$ is strictly increasing.

  Now consider $Q(s) = \frac{H(s)}{c}$. This $H$ is
  obviously also differentiable, since $H$ is. The equation
  for $H$ yields the following equation for $Q$,
  \begin{align*}
    Q(h+s) = Q(h) \alpha(s) + Q(s) \qquad \text{ for any } s,h \in \bR ,
  \end{align*}
  where $\alpha(s) = A(e^{s})$.
  By Lemma~\ref{lem:ode-of-order-two-for-Q} $Q$ is then twice
  continuously differentiable and satisfies
  \begin{align}
  Q'(s) Q'' (0) = Q'' (s) \qquad \text{ for all } s \in \bR .
  \end{align}
  By Lemma~\ref{lem:solve-Q} we then get
  a formula for $Q$, which involves
  \begin{align*}
  \gamma = Q''(0) = \frac{H''(0)}{c}
    = \frac{1}{c} \dder{s} E(e^s) \big|_{s = 0} .
  \end{align*}

  If $\gamma \ne 0$ then the formula reads
  $Q(s) = \frac{e^{\gamma s} - 1}{\gamma}$.
  Now just tracing the definitions, we get the asserted formula
  \begin{align*}
  E(\lambda) = H(\log \lambda)
    = c \, Q(\log \lambda) = c \, \big( e^{\gamma \, \log(\lambda)} - 1 \big)
    = c \, \big( \lambda^{\gamma} - 1 \big) .
  \end{align*}

  If $\gamma = 0$ then the formula reads $Q(s) = s$.
  Now just tracing the definitions, we get instead
  \begin{align*}
  E(\lambda) = H(\log \lambda)
    = c \, Q(\log \lambda) = c \, \log(\lambda) ,
  \end{align*}
  and the proof is complete.
\end{proof}


\section{Inverting to recover a cumulative distribution function}

%\begin{lemma}[Classification of cdfs of evds]
%  \label{lem:cdf-of-evd-classification}
%  %\uses{}
%  %\lean{}
%  % \leanok
%  Suppose that $G \colon \bR \to [0,1]$ is a c.d.f.
%  whose right-continuous inverse is given by $G^\leftarrow = \cdots$
%\end{lemma}


%%FILE blueprint/src/sections/extreme_value_distributions.tex
\chapter{Extreme value distributions}

\section{Definition of extreme value distributions}

Lemma~\ref{lem:cdf-of-affine-max-many} motivates the following definition.

\begin{definition}
  \label{def:extr-val-distr}
  \uses{lem:oriented-affine-action-on-cdf}
  \lean{CumulativeDistributionFunction.IsExtremeValueDistr}
  \leanok
  A c.d.f. $G$ is said to be an \emph{extreme value distribution} if
  $G$ is nondegenerate and
  there exists a c.d.f. $F$ and a sequence $(A_n)_{n \in \N}$ of orientation
  preserving affine isomorphisms $A_n \in \OriAffR$, such that
  for every continuity point $x \in \R$ of~$G$ we have
  \begin{align*}
  \lim_{n \to \infty} \big( (A_n.F)(x) \big)^n \; = \; G(x) .
  \end{align*}
\end{definition}

\begin{lemma}
  \label{lem:extr-val-distr-transform}
  \uses{lem:oriented-affine-action-on-cdf, def:extr-val-distr}
  \lean{CumulativeDistributionFunction.IsExtremeValueDistr.affineTransform}
  \leanok
  Let $G$ be an extreme value distribution and
  and $A \in \OriAffR$ an orientation preserving affine isomorphism.
  Then also $A.G$ is an extreme value distribution.
\end{lemma}
\begin{proof}
  \uses{lem:cdf-continuity-pt-transform, lem:degenerate-cdf-transform}
  \leanok
  Straightforward using Lemmas~\ref{lem:cdf-continuity-pt-transform}
  and~\ref{lem:degenerate-cdf-transform}.
\end{proof}



\section{Three types of extreme value distributions}

\begin{definition}
  \label{def:std-Gumbel-cdf}
  \uses{def:cdf}
  \lean{standardGumbelCDF}
  % \leanok
  \emph{The standard Gumbel distribution} is the c.d.f. $\GumbelCDF$ given by
  \begin{align*}
  \GumbelCDF (x) = \exp\big(-\exp(-x)\big) .
  \end{align*}

  (In the parametrization of extreme value distribution types by one
  index $\gamma \in \bR$, this case corresponds to $\gamma = 0$.)
\end{definition}

\begin{definition}
  \label{def:std-Weibull-cdf}
  \uses{def:cdf}
  \lean{standardWeibullCDF}
  % \leanok
  \emph{The standard (reverse) Weibull distribution} of parameter
  $\alpha > 0$ is the c.d.f. $\WeibullCDF{\alpha}$ given by
  \begin{align*}
  \WeibullCDF{\alpha} (x) = \begin{cases}
      \exp \big( - (-x)^\alpha \big) \quad & \text{ for } x < 0 \\
      1 \quad & \text{ for } x \ge 0 .
    \end{cases}
  \end{align*}

  (In the parametrization of extreme value distribution types by one
  index $\gamma \in \bR$, this case corresponds to $\gamma < 0$
  via $\gamma = -1/\alpha$.)
\end{definition}

\begin{definition}
  \label{def:std-Frechet-cdf}
  \uses{def:cdf}
  \lean{standardFrechetCDF}
  % \leanok
  \emph{The standard Fr\'echet distribution} of parameter $\alpha > 0$
  is the c.d.f. $\FrechetCDF{\alpha}$ given by
  \begin{align*}
  \FrechetCDF{\alpha} (x) = \begin{cases}
    0 \quad & \text{ for } x \le 0 \\
    \exp \big( - x^{-\alpha} \big) \quad & \text{ for } x > 0 .
    \end{cases}
  \end{align*}

  (In the parametrization of extreme value distribution types by one
  index $\gamma \in \bR$, this case corresponds to $\gamma > 0$
  via $\gamma = 1/\alpha$.)
\end{definition}

\begin{theorem}
  \label{thm:Gumbel-is-extr-val-distr}
  \uses{def:std-Gumbel-cdf, def:extr-val-distr}
  \lean{isExtremeValueDistr_standardGumbelCDF}
  \leanok
  The standard Gumbel distribution $\GumbelCDF$ is an extreme value
  distribution.
\end{theorem}
\begin{proof}
  %\uses{lem:cdf-of-max-many}
  % \leanok
  Set $A_n(x) = x - \log(n)$ for $n \in \bN$. Then
  $A_n^{-1}(x) = x + \log(n)$ and for any $n \ge 1$ and $x \in \bR$ we get
  \begin{align*}
    \big( (A_n . \GumbelCDF)(x) \big)^n
    = \; & \big( \GumbelCDF(x + \log(n)) \big)^n \\
    = \; & \Big( \exp \big( - \exp(-(x + \log n)) \big) \Big)^n \\
    = \; & \exp \big( - n \exp(-x - \log n) \big) \\
    = \; & \exp \big( - n \, e^{-x} \, e^{-\log n}  \big) \\
    %= \; & \exp \big( - n \, e^{-x} \, n^{-1} \big) \\
    = \; & \exp \big( - e^{-x} \big) \\
    = \; & \GumbelCDF(x) .
  \end{align*}
  Since the above is true for each $n$, we in particular have
  \begin{align*}
    \lim_{n \to \infty} \big( (A_n . \GumbelCDF)(x) \big)^n
    = \GumbelCDF(x)
  \end{align*}
  for all $x \in \bR$. Since $\GumbelCDF$ is also nondegenerate,
  this shows that it is an extreme value distribution.
\end{proof}

\begin{theorem}
  \label{thm:Weibull-is-extr-val-distr}
  \uses{def:std-Weibull-cdf, def:extr-val-distr}
  \lean{isExtremeValueDistr_standardWeibullCDF}
  % \leanok
  For any $\alpha > 0$, the standard Weibull distribution $\WeibullCDF{\alpha}$
  is an extreme value distribution.
\end{theorem}
\begin{proof}
  %\uses{lem:cdf-of-max-many}
  % \leanok
  \ldots
\end{proof}

\begin{theorem}
  \label{thm:Frechet-is-extr-val-distr}
  \uses{def:std-Frechet-cdf, def:extr-val-distr}
  \lean{isExtremeValueDistr_standardFrechetCDF}
  % \leanok
  For any $\alpha > 0$, the standard Fr\'echet distribution $\FrechetCDF{\alpha}$
  is an extreme value distribution.
\end{theorem}
\begin{proof}
  \uses{}
  % \leanok
  \ldots
\end{proof}



\section{Equivalent formulations of the limit relation}

\begin{lemma}[Logarithmic version of the limit relation]
  \label{lem:log-ev-limit}
  \uses{def:cdf, lem:oriented-affine-action-on-cdf}
  \lean{ev_limit_iff_log_ev_limit}
  \leanok
  Let $F$ and $G$ be c.d.f.s, and $(A_n)_{n \in \N}$ a sequence of orientation
  preserving affine isomorphisms $A_n \in \OriAffR$. Then for any $x \in \bR$
  such that $0 < G(x) < 1$, the two conditions
  \begin{align*}
  \text{(i)} & & \lim_{n \to \infty} \big( (A_n.F)(x) \big)^n \, = \; & G(x) \\
  \text{(ii)} & & \lim_{n \to \infty} n \, \log F (A_n^{-1}(x)) \, = \; & \log G(x)
  \end{align*}
  are equivalent.
\end{lemma}
\begin{proof}
  \uses{}
  \leanok
  Recall that $(A_n.F)(x) = F (A_n^{-1}(x))$.
  Then just take logarithms (and use continuity) to get from (i) to (ii),
  and take exponentials (and use continuity) to get from (ii) to (i).
\end{proof}

\begin{lemma}[Relation implies $F$ tending to one]
  \label{lem:ev-limit-cdf-affine-tendsto-one}
  \uses{def:cdf, lem:oriented-affine-action-on-cdf}
  \lean{tendsto_one_of_ev_limit}
  \leanok
  Let $F$ and $G$ be c.d.f.s, and $(A_n)_{n \in \N}$ a sequence of orientation
  preserving affine isomorphisms $A_n \in \OriAffR$. Then for any $x \in \bR$
  such that $0 < G(x) < 1$, if
  \begin{align*}
  \text{(i)} & & \lim_{n \to \infty} \big( (A_n.F)(x) \big)^n \, = \; & G(x)
  \end{align*}
  holds, then necessarily
  \begin{align*}
    \lim_{n \to \infty} F (A_n^{-1}(x)) \; = \; 1 .
  \end{align*}
\end{lemma}
\begin{proof}
  % \uses{}
  % \leanok
  Otherwise $\big( F (A_n^{-1}(x)) \big)^n$ would have $0 \ne G(x)$
  as an accumulation point, contradicting the assumed limit (i).

  To wit, if for some $\delta > 0$ we would have
  $F (A_n^{-1}(x)) \le 1 - \delta$ for infinitely many $n$, then
  $0 \le \big( F (A_n^{-1}(x)) \big)^n \le (1 - \delta)^n $
  for those $n$, and since $(1 - \delta)^n \to 0$, we would
  get $\big( F (A_n^{-1}(x)) \big)^n \to 0 \ne G(x)$ along the
  subsequence of those $n$; a contradiction.
\end{proof}

\begin{lemma}[Taylor expansion limit modification]
  \label{lem:modify-limit-taylor}
  \lean{tendsto_smul_apply_smul_deriv_of_tendsto_atTop_of_tendsto_smul_apply_smul_deriv,
  tendsto_zero_of_tendsto_atTop_of_tendsto_smul}
  \leanok
  Let $S \subset \bR$ be a subset with $0 \in S$, and let $f_1, f_2 \colon S \to \bR$ be
  functions. Let also $(t_n)_{n \in \bN}$ be a sequence in $S$,
  and let $(m_n)_{n \in \bR}$ be a sequence of real numbers tending to infinity,
  $\lim_{n \to \infty} m_n = +\infty$.
  Assume that for $j = 1,2$
  \begin{itemize}
    \item $f_j(0)=0$;
    \item the derivative $f_j'(0)$ exists (derivative taken within the set $S$)
  \end{itemize}
  Assume further about $j = 1$ that
  \begin{itemize}
    \item $f_1'(0) \ne 0$;
    \item the limit $\lim_{n \to \infty} \big( m_n \, f_1 (t_n) \big)$ exists;
    \item for any $\delta > 0$ there exists an $\eps > 0$ such that
      if $|f_1(t)| < \eps$ (with $t \in S$) then $|t| < \delta$.
  \end{itemize}
  Denote $c = \frac{1}{f_1'(0)} \, \lim_{n \to \infty} \big( m_n \, f_1 (t_n) \big)$.
  Then we have $\lim_{n \to 0} t_n = 0$ and
  \begin{align*}
    \lim_{n \to \infty} \Big( m_n \, f_2 (t_n) \Big) \; = \; & c \, f_2'(0) .
  \end{align*}
\end{lemma}
\begin{proof}
  % \uses{}
  \leanok
  This is in principle straightforward:
  the assumptions are first checked to imply that $\lim_{n \to 0} t_n = 0$,
  and then one can just consider the first order Taylor expansions of the
  functions $f_j$ at $0$ given by the assumed existence of the derivatives
  (the key is $t_n \, = \, \frac{c}{m_n} + \oo \big(\frac{1}{m_n}\big)$).
\end{proof}

\begin{lemma}[Taylored version of the limit relation]
  \label{lem:taylored-ev-limit}
  \uses{def:cdf, lem:oriented-affine-action-on-cdf}
  \lean{log_ev_limit_iff_taylored_ev_limit}
  \leanok
  Let $F$ and $G$ be c.d.f.s, and $(A_n)_{n \in \N}$ a sequence of orientation
  preserving affine isomorphisms $A_n \in \OriAffR$. Then for any $x \in \bR$
  such that $0 < G(x) < 1$, the two conditions
  \begin{align*}
    \text{(ii)} & & \lim_{n \to \infty} n \, \log F (A_n^{-1}(x)) \, = \; & \log G(x) \\
    \text{(iii)} & & \lim_{n \to \infty} \Big( n \, \big(1 - F (A_n^{-1}(x)) \big) \Big)
      \, = \; & - \log G(x)
  \end{align*}
  are equivalent.
\end{lemma}
\begin{proof}
  \uses{lem:modify-limit-taylor}
  % \leanok

  Both implications {(ii) $\, \Rightarrow \,$ (iii)} and {(iii) $\, \Rightarrow \,$ (ii)}
  are proven similarly using Lemma~\ref{lem:modify-limit-taylor}.

  Assume (ii). Let $f_1(t) = - \log (1-t)$ and $f_2(t) = t$ and
  $S = [0,1) \subset \bR$, and $m_n = n$ and $t_n = 1 - F (A_n^{-1}(x))$.
  It is straightforward to check the assumptions of Lemma~\ref{lem:modify-limit-taylor},
  with
  \begin{align*}
    f_1'(0) \; = \; & \der{t} \Big|_{t=0} \Big( - \log (1 - t) \Big) \; = \; 1
  \end{align*}
  and $f_2'(0) = \mathrm{id}'(0) = 1$.
  The key assumption about the existence of
  the limit $\lim_{n \to \infty} \big( m_n \, f_1(t_n) \big)$ is given by (ii),
  and the conclusion is (iii).

  Similarly assuming (iii) we derive (ii) with Lemma~\ref{lem:modify-limit-taylor}
  just interchanging the roles of the two functions, i.e., now using
  $f_1(t) = t$ and $f_2(t) = - \log (1-t)$ instead.
\end{proof}

\begin{lemma}[Inverted Taylored version of the limit relation]
  \label{lem:inv-taylored-ev-limit}
  \uses{def:cdf, lem:oriented-affine-action-on-cdf}
  \lean{taylored_ev_limit_iff_oneDivOneSub_limit}
  \leanok
  Let $F$ and $G$ be c.d.f.s, and $(A_n)_{n \in \N}$ a sequence of orientation
  preserving affine isomorphisms $A_n \in \OriAffR$. Then for any $x \in \bR$
  such that $0 < G(x) < 1$, the two conditions
  \begin{align*}
    \text{(iii)} & & \lim_{n \to \infty} \Big( n \, \big(1 - F (A_n^{-1}(x)) \big) \Big)
      \, = \; & - \log G(x) \\
    \text{(iv)} & & \lim_{n \to \infty} \frac{1}{n \, \big(1 - F (A_n^{-1}(x)) \big)}
      \, = \; & \frac{1}{-\log G(x)}
  \end{align*}
  are equivalent.
\end{lemma}
\begin{proof}
  % \uses{}
  \leanok
  By assumption $G(x) \in (0,1)$ we have $-\log G(x) > 0$.

  The implication {(iii)~$\,\Rightarrow \,$~(iv)}
  can therefore be seen by applying $t \mapsto \frac{1}{t}$
  and using its continuity at $t = -\log G(x)$, and the
  converse implication {(iv)~$\,\Rightarrow \,$~(iii)} similarly by using continuity of
  $t \mapsto \frac{1}{t}$ at $t = \frac{1}{-\log G(x)}$.
\end{proof}

\begin{lemma}[Transformed version of the limit relation]
  \label{lem:transformed-ev-limit}
  \uses{def:one-div-one-sub-cdf, def:one-div-neg-log-cdf, lem:oriented-affine-action-on-cdf}
  \lean{oneDivSub_limit_iff}
  \leanok
  Let $F$ and $G$ be c.d.f.s, and $(A_n)_{n \in \N}$ a sequence of orientation
  preserving affine isomorphisms $A_n \in \OriAffR$. Then for any $x \in \bR$
  such that $0 < G(x) < 1$, the two conditions
  \begin{align*}
    \text{(iv)} & & \lim_{n \to \infty} \frac{1}{n \, \big(1 - F (A_n^{-1}(x)) \big)}
      \, = \; & \frac{1}{-\log G(x)} \\
    \text{(v)} & & \lim_{n \to \infty} \frac{1}{n} \oneDivOneSub{A_n . F} (x)
      \, = \; & \oneDivNegLog{G} (x)
  \end{align*}
  are equivalent.

  (See Definitions~\ref{def:one-div-one-sub-cdf}
  and \ref{def:one-div-neg-log-cdf} for the transforms involved in condition (v)).
\end{lemma}
\begin{proof}
  % \uses{}
  \leanok
  This is in principle straightforward, although certain cases need to be checked
  separately and the continuity of various natural extensions need to addressed.
\end{proof}

% def:one-div-neg-log-cdf


\begin{theorem}[Equivalent versions of the limit relation]
  \label{thm:tfae-ev-limit}
  \uses{def:one-div-one-sub-cdf, def:one-div-neg-log-cdf, lem:oriented-affine-action-on-cdf}
  \lean{tfae_ev_limit_relation}
  \leanok
  Let $F$ and $G$ be c.d.f.s, and $(A_n)_{n \in \N}$ a sequence of orientation
  preserving affine isomorphisms $A_n \in \OriAffR$. Then for any $x \in \bR$
  such that $0 < G(x) < 1$, the conditions
  \begin{align*}
    \text{(i)} & & \lim_{n \to \infty} \big( (A_n.F)(x) \big)^n \, = \; & G(x) \\
    \text{(ii)} & & \lim_{n \to \infty} n \, \log F (A_n^{-1}(x)) \, = \; & \log G(x) \\
    \text{(iii)} & & \lim_{n \to \infty} \Big( n \, \big(1 - F (A_n^{-1}(x)) \big) \Big)
      \, = \; & - \log G(x) \\
    \text{(iv)} & & \lim_{n \to \infty} \frac{1}{n \, \big(1 - F (A_n^{-1}(x)) \big)}
      \, = \; & \frac{1}{-\log G(x)} \\
    \text{(v)} & & \lim_{n \to \infty} \frac{1}{n} \oneDivOneSub{A_n . F} (x)
      \, = \; & \oneDivNegLog{G} (x)
  \end{align*}
  are equivalent.

  (See Definitions~\ref{def:one-div-one-sub-cdf}
  and \ref{def:one-div-neg-log-cdf} for the transforms involved in condition (v)).

  \emph{(More equivalent conditions are to be added; this is just a theorem to collect
  various equivalent phrasings.)}
\end{theorem}
\begin{proof}
  \uses{lem:taylored-ev-limit, lem:log-ev-limit, lem:inv-taylored-ev-limit,
  lem:transformed-ev-limit}
  \leanok
  This is just a combination of earlier results.
\end{proof}


%%FILE blueprint/src/sections/lcinverse.tex
\chapter{Left-continuous inverses}

Let $R$ and $S$ be complete linear orders, for example $[0,1]$, $[0,+\infty]$, or
$[-\infty,+\infty] =: \overline{\bR}$.

\begin{definition}
  \label{def:lc-inverse}
  \lean{lcInv}
  \leanok
  Let $f \colon R \to S$ be a function (usually assumed nondecreasing).
  The left-continuous inverse of $f$ is the function $\lcinv{f} \colon S \to R$
  given by
  \begin{align*}
    \lcinv{f}(y) := \inf \big\{ x \in R \; \big| \; f(x) \ge y \big\}
    \qquad \text{ for } y \in S.
  \end{align*}
  The right-continuous inverse $\rcinv{f} \colon S \to R$ is analoguously defined by
  \begin{align*}
    \rcinv{f}(y) := \sup \big\{ x \in R \; \big| \; f(x) \le y \big\}
    \qquad \text{ for } y \in S.
  \end{align*}
\end{definition}


%%FILE blueprint/src/sections/one_parameter_affine.tex
\chapter{One-parameter subgroups of affine isomorphisms}

%The group $\OriAffR$ is isomorphic to the matrix group
%\begin{align*}
%  \left\{ \left[ \begin{array}{cc} a & b \\ 0 & 1 \end{array} \right] \; \Bigg|
%    \; a>0, b \in \bR \right\} .
%\end{align*}
%This is a 2-dimensional Lie group with Lie algebra\ldots


\begin{lemma}[Functional equation in one parameter subgroups of affine isomorphisms]
  \label{lem:affine-one-parameter-subgroup-functional-eqn}
  \uses{def:oriented-affine-isomorphism}
  % \lean{}
  % \leanok
  Suppose that $t \mapsto A_t$ is a homomorphism of multiplicative groups
  $(0,+\infty) \to \OriAffR$, i.e., for any $s, t > 0$ we have
  \begin{align*}
    A_{s t} = A_s \circ A_t .
  \end{align*}
  Write $A_t(x) = a_t x + b_t$, with $a_t > 0$ and $b_t \in \bR$.
  Then we have, for any $s, t > 0$,
  \begin{align*}
    a_{t s} = \; & a_t \, a_s
      \qquad \text{ and } \\
    b_{t s} = \; & a_t \, b_s + b_t .
  \end{align*}
  (Also by symmetry $b_{t s} = a_s \, b_t + b_s$.)
\end{lemma}
\begin{proof}
  % \uses{}
  % \leanok
\end{proof}

\begin{lemma}[Functional equation scaling coefficient solution]
  \label{lem:solution-functional-eqn-scaling}
  %\uses{def:lem:oriented-affine-action-on-cdf}
  \lean{eq_exp_const_mul_of_multiplicative_of_measurable}
  \leanok
  Suppose that $a \colon (0,+\infty) \to (0,+\infty)$ is a
  measurable function satisfying, for any $s, t > 0$,
  \begin{align*}
    a(t s) = \; & a(t) \, a(s) .
  \end{align*}
  Then there exists a $\rho \in \bR$ such that for all $t > 0$,
  \begin{align*}
    a(t) = t^{\rho} .
  \end{align*}
\end{lemma}
\begin{proof}
  \uses{lem:different-difference-set-contains-interval}
  % \leanok
\end{proof}

\begin{lemma}[Functional equation translation coefficient solution with $\rho = 0$]
  \label{lem:solution-functional-eqn-translation-no-scaling}
  %\uses{def:lem:oriented-affine-action-on-cdf}
  \lean{eq_const_mul_of_additive_of_measurable,
        eq_exp_const_mul_of_multiplicative_of_measurable}
  \leanok
  Suppose that
  $b \colon (0,+\infty) \to \bR$ is a measurable function
  satisfying, for any $s, t > 0$,
  \begin{align*}
    b(t s) = \; & b(s) + b(t) .
  \end{align*}
  Then there exists a constant $c$ such that for $t \in (0,+\infty)$
  we have
  \begin{align*}
    b(t) = \; & -c \, \log (t) .
  \end{align*}
\end{lemma}
\begin{proof}
  \uses{lem:solution-functional-eqn-scaling}
  \leanok
\end{proof}

\begin{lemma}[Functional equation translation coefficient solution with $\rho \ne 0$]
  \label{lem:solution-functional-eqn-translation-with-scaling}
  %\uses{def:lem:oriented-affine-action-on-cdf}
  \lean{eq_const_mul_one_sub_exp_of_multiplicative_with_scaling_of_measurable}
  \leanok
  Suppose that $\rho \in \bR \setminus \set{0}$ and
  $b \colon (0,+\infty) \to \bR$ is a measurable function
  satisfying, for any $s, t > 0$,
  \begin{align*}
    b(t s) = \; & t^{\rho} \, b(s) + b(t) .
  \end{align*}
  Then there exists a constant $c$ such that for $t \in (0,+\infty) \setminus \set{1}$
  we have
  \begin{align*}
    b(t) = \; & c (1 - t^{-\rho}) .
  \end{align*}
\end{lemma}
\begin{proof}
  \uses{lem:solution-functional-eqn-scaling}
  \leanok
\end{proof}

The above solutions to functional equations can be used to classify all
one-parameter subgroups of the group of oriented affine isomorphisms of $\bR$.
Such a subgroup can be given in a parametrized form as a group homomorphism
$\bR \to \OriAffR$ (from the additive group $\bR$) or alternatively
as group homomorphisms $(0,+\infty) \to \OriAffR$ (from the multiplicative
group $(0,+\infty)$). The additive and multiplicative versions are
related by the change of variable
$\bR \ni t \leftrightarrow \lambda := e^t \in (0,+\infty)$
(conversely, $t = \log(\lambda)$).
\emph{(In Lean the type $\bR$ is more convenient than the type $(0,+\infty)$,
so in formal statements we prefer the former choice.)}

\begin{theorem}[One-parameter subgroups of affine isomorphisms of $\bR$]
  \label{thm:one-parameter-subgroups-of-affine-isomorphisms}
  \uses{def:oriented-affine-isomorphism, def:translation-subgroup, def:fixing-subgroup}
  \lean{AffineIncrEquiv.homomorphism_from_Real_characterization,
    AffineIncrEquiv.homomorphism_from_Real_characterization_of_nontrivial}
  \leanok
  \emph{[TODO: Switch to additive notation and $\bR$ rather than multiplicative
  notation and $(0,+\infty)$, to match the most convenient formalized statements.]}

  Suppose that $t \mapsto A_t$ is a measurable homomorphism of multiplicative groups
  $(0,+\infty) \to \OriAffR$, i.e., for any $s, t > 0$ we have
  \begin{align*}
    A_{s t} = A_s \circ A_t
  \end{align*}
  and $A_t(x) = a_t x + b_t$, with $t \mapsto a_t$ and $t \mapsto b_t$
  measurable functions.
  Then either
  \begin{description}
    \item[(0)] there exists a $\beta \in \bR$ such that for all $t > 0$ and $x \in \bR$
      \begin{align*}
        A_t(x) = x + \beta \, \log(t) ;
      \end{align*}
      or
    \item[(1)] there exists a $\rho \ne 0$ and $c \in \bR$ such that
      for all $t > 0$ and $x \in \bR$
      \begin{align*}
        A_t(x) = t^{\rho} (x - c) + c .
      \end{align*}
  \end{description}
\end{theorem}
\begin{proof}
  \uses{lem:affine-one-parameter-subgroup-functional-eqn,
    lem:solution-functional-eqn-scaling,
    lem:solution-functional-eqn-translation-no-scaling,
    lem:solution-functional-eqn-translation-with-scaling}
  \leanok
\end{proof}


%%FILE blueprint/src/sections/transforms.tex
\chapter{Transforms of cumulative distribution functions}

In this part, we introduce certain transforms and extensions of
cumulative distribution functions, which are used in the classification
calculation of the extreme value distributions.


\section{Extended cumulative distribution function}

\begin{definition}[Extended cumulative distribution function]
  \label{def:cdf-extend}
  \uses{def:cdf}
  \lean{CumulativeDistributionFunction.extend}
  \leanok
  The extension $\extendcdf{F}$ of a c.d.f. $F$ is the function
  \begin{align*}
    \extendcdf{F} \colon [-\infty,+\infty] \to [0,1]
  \end{align*}
  given by
  \begin{align*}
    \extendcdf{F}(x) = \begin{cases}
      0 & \; \text{ if } x = -\infty \\
      F(x) & \; \text{ if } -\infty < x < +\infty \\
      1 & \; \text{ if } x = +\infty .
      \end{cases}
  \end{align*}
\end{definition}

\begin{lemma}[Continuity points of extended c.d.f.]
  \label{lem:cdf-extend-continuity-pts}
  \uses{def:cdf-extend}
  \lean{CumulativeDistributionFunction.extend_continuousAt,
  CumulativeDistributionFunction.extend_continuousAt_top,
  CumulativeDistributionFunction.extend_continuousAt_bot}
  \leanok
  The extension $\extendcdf{F}$ of a c.d.f. $F$ is continuous
  at $x = -\infty$, $x = + \infty$, and at any
  $x \in (-\infty,+\infty)$ where $F$ is continuous.
\end{lemma}
\begin{proof}
  % \uses{}
  % \leanok
  Since $\lim_{x \to +\infty} F(x) = 1$ by properties of c.d.f.s
  and $\extendcdf{F}(+\infty) = 1$ by definition of the extension,
  continuity at $x = + \infty$ follows. Continuity at $x=-\infty$
  is similar.

  Suppose $F$ is continuous at $x \in \bR$.
  %Let $\iota \colon \bR \to [-\infty,+\infty]$ be the inclusion
  %of the real line into the extended real line.
  Then since $\extendcdf{F}$ coincides with $F$
  in a neighborhood of $x$
  (indeed on the open set $\bR \subsetneq [-\infty,+\infty]$),
  the continuity of $F$ at $x \in \bR$ implies
  continuity of $\extendcdf{F}$ at $x$.
\end{proof}

\section{One over one minus cumulative distribution function}

\begin{definition}[One over one minus cumulative distribution function]
  \label{def:one-div-one-sub-cdf}
  \uses{def:cdf-extend}
  \lean{CumulativeDistributionFunction.oneDivOneSub}
  \leanok
  The transform $\oneDivOneSub{F}$ of a c.d.f. $F$ is the function
  \begin{align*}
    \oneDivOneSub{F} \colon [-\infty,+\infty] \to [1,+\infty]
  \end{align*}
  given by
  \begin{align*}
    \oneDivOneSub{F}(x) = \begin{cases}
      1 & \; \text{ if } \extendcdf{F}(x) = 0 \\
      \frac{1}{1-\extendcdf{F}(x)} & \; \text{ if } 0 < \extendcdf{F}(x) < 1 \\
      +\infty & \; \text{ if } \extendcdf{F}(x) = 1 ,
      \end{cases}
  \end{align*}
  where $\extendcdf{F} \colon [-\infty,+\infty] \to [0,1]$ is the extension of the c.d.f. $F$.
\end{definition}

\begin{lemma}[Continuity points of one over one minus c.d.f.]
  \label{lem:one-div-one-sub-cdf-continuity-pts}
  \uses{def:one-div-one-sub-cdf}
  \lean{CumulativeDistributionFunction.oneDivOneSub_continuousAt,
  CumulativeDistributionFunction.oneDivOneSub_continuousAt_top,
  CumulativeDistributionFunction.oneDivOneSub_continuousAt_bot}
  \leanok
  The transform $\oneDivOneSub{F}$ of a c.d.f. $F$ is continuous
  at $x = -\infty$, $x = + \infty$, and at any
  $x \in (-\infty,+\infty)$ where $F$ is continuous.
\end{lemma}
\begin{proof}
  \uses{lem:cdf-extend-continuity-pts}
  % \leanok
  Since $\lim_{x \to +\infty} \extendcdf{F}(x) = \extendcdf{F}(+\infty) = 1$
  by Lemma~\ref{lem:cdf-extend-continuity-pts} and the continuous extension
  of $p \mapsto \frac{1}{1-p}$ to a function $[0,1] \to [0,+\infty]$
  tends to $+\infty$ as $p \to 1$, we have
  \begin{align*}
    \lim_{x \to +\infty} \oneDivOneSub{F}(x)
    \; = \; +\infty \; = \; \oneDivOneSub{F}(+\infty) .
  \end{align*}
  Therefore $\oneDivOneSub{F}$ is continuous at $+\infty$.
  Continuity at $-\infty$ similarly follows from
  $\lim_{x \to -\infty} \extendcdf{F}(x) = \extendcdf{F}(-\infty) = 0$
  and $\frac{1}{1-p}$ tending to $1$
  as $p \to 0$, which give
  \begin{align*}
    \lim_{x \to -\infty} \oneDivOneSub{F}(x)
    \; = \; 1 \; = \; \oneDivOneSub{F}(-\infty) .
  \end{align*}

  Suppose $F$ is continuous at $x \in \bR$,
  and recall from Lemma~\ref{lem:cdf-extend-continuity-pts} that
  $\extendcdf{F}$ is then also continuous at $x$.
  Now $\oneDivOneSub{F}$ is a composition of the continuous function
  $p \mapsto \frac{1}{1-p} \colon [0,1] \to [0,+\infty]$ with
  $\extendcdf{F}$, and as such also becomes continuous at $x$.
\end{proof}

\section{One over negative logarithm cumulative distribution function}

\begin{definition}[One over negative logarithm cumulative distribution function]
  \label{def:one-div-neg-log-cdf}
  \uses{def:cdf-extend}
  \lean{CumulativeDistributionFunction.oneDivNegLog}
  \leanok
  The transform $\oneDivNegLog{F}$ of a c.d.f. $F$ is the function
  \begin{align*}
    \oneDivNegLog{F} \colon [-\infty,+\infty] \to [0,+\infty]
  \end{align*}
  given by
  \begin{align*}
    \oneDivNegLog{F}(x) = \begin{cases}
      0 & \; \text{ if } \extendcdf{F}(x) = 0 \\
      \frac{1}{\log \big( 1/\extendcdf{F}(x) \big)} & \; \text{ if } 0 < \extendcdf{F}(x) < 1 \\
      +\infty & \; \text{ if } \extendcdf{F}(x) = 1 ,
      \end{cases}
  \end{align*}
  where $\extendcdf{F} \colon [-\infty,+\infty] \to [0,1]$ is the extension of the c.d.f. $F$.
\end{definition}

\begin{lemma}[Continuity points of one over negative logarithm c.d.f.]
  \label{lem:one-div-neg-log-cdf-continuity-pts}
  \uses{def:one-div-neg-log-cdf}
  \lean{CumulativeDistributionFunction.oneDivNegLog_continuousAt,
  CumulativeDistributionFunction.oneDivNegLog_continuousAt_top,
  CumulativeDistributionFunction.oneDivNegLog_continuousAt_bot}
  \leanok
  The transform $\oneDivNegLog{F}$ of a c.d.f. $F$ is continuous
  at $x = -\infty$, $x = + \infty$, and at any
  $x \in (-\infty,+\infty)$ where $F$ is continuous.
\end{lemma}
\begin{proof}
  \uses{lem:cdf-extend-continuity-pts}
  % \leanok
  \uses{lem:cdf-extend-continuity-pts}
  % \leanok
  Since $\lim_{x \to +\infty} \extendcdf{F}(x) = \extendcdf{F}(+\infty) = 1$
  by Lemma~\ref{lem:cdf-extend-continuity-pts} and the continuous extension
  of $p \mapsto \frac{1}{\log(1/p)}$ to a function $[0,1] \to [0,+\infty]$
  tends to $+\infty$ as $p \to 1$, we have
  \begin{align*}
    \lim_{x \to +\infty} \oneDivNegLog{F}(x)
    \; = \; +\infty \; = \; \oneDivNegLog{F}(+\infty) .
  \end{align*}
  Therefore $\oneDivNegLog{F}$ is continuous at $+\infty$.
  Continuity at $-\infty$ similarly follows from
  $\lim_{x \to -\infty} \extendcdf{F}(x) = \extendcdf{F}(-\infty) = 0$
  and $\frac{1}{\log(1/p)}$ tending to $0$
  as $p \to 0$, which give
  \begin{align*}
    \lim_{x \to -\infty} \oneDivNegLog{F}(x)
    \; = \; 0 \; = \; \oneDivNegLog{F}(-\infty) .
  \end{align*}

  Suppose $F$ is continuous at $x \in \bR$,
  and recall from Lemma~\ref{lem:cdf-extend-continuity-pts} that
  $\extendcdf{F}$ is then also continuous at $x$.
  Now $\oneDivNegLog{F}$ is a composition of the continuous function
  $p \mapsto \frac{1}{\log(1/p)} \colon [0,1] \to [0,+\infty]$ with
  $\extendcdf{F}$, and as such also becomes continuous at $x$.
\end{proof}


%%FILE blueprint/src/sections/types.tex
\chapter{Types}

\begin{definition}[Type (of distribution on $\bR$)]
  \label{def:cdf-type}
  \uses{def:cdf, lem:oriented-affine-action-on-cdf}
  %\lean{lcinv}
  %\leanok
  Two c.d.f.s $F, G$ are said to be of the same type, if
  there exists an order-preserving affine isomorphism $A \in \OriAffR$
  such that $G = A . F$.

  Being of the same type is an equivalence relation, and
  the equivalence classes are called types (of distributions on $\bR$).
\end{definition}


\section{Convergence to types}

The notion of convergence of cumulative distribution functions
considered here is always taken to be pointwise convergence
on the set of continuity points of the limit c.d.f.
By Theorem~\ref{thm:convergence-in-distribution-with-cdf} and
Lemma~\ref{lem:cdf-convergence-from-convergence-in-distribution},
this corresponds to convergence in distribution (weak convergence
of probability measures). Below, when we write $F_n \tendsinlaw F$
for c.d.f.s $F_n$, $n \in \bN$, and $F$, this is always what we
mean.

For affine maps $A \colon \bR \to \bR$, we use the topology
of pointwise convergence of functions. Equivalently, convergence of affine maps
means the convergence of the coefficients $a, b \in \bR$ in the expression
$x \mapsto a x + b$ of the affine functions
(so $A_n \to A$ if and only if the functions are of the form
$A_n(x) = a_n x + b_n$ and $A(x) = a x + b$, and $a_n \to a$ and $b_n \to b$).

\begin{lemma}[Unique affine relation among two nondegenerate c.d.f.s]
  \label{lem:unique-affine-relation-cdf}
  \uses{def:cdf, def:degenerate-cdf, lem:oriented-affine-action-on-cdf}
  \lean{CumulativeDistributionFunction.unique_orientationPreservingAffineEquiv_smul_eq_not_isDegenerate}
  \leanok
  Let $F, G$ be two c.d.f.s of the same type, and $A \in \OriAffR$
  an affine isomorphism such that $G = A.F$. If $F$ is nondegenerate,
  then $A$ is the only element of $\OriAffR$ for which the
  relation $G = A.F$ holds.
\end{lemma}
\begin{proof}
  % \uses{}
  % \leanok
  Since $F$ is nondegenerate, we can find two different points $x_1 < x_2$
  such that $0 < F(x_1) < F(x_2) \le 1$. By right continuity of $F$, we can
  assume these points to be taken minimal with the given values, i.e.,
  $x_j = \inf \big\{ x \in \bR \; \big| \; F(x) = F(x_j) \big\}$ for
  $j=1,2$.

  The assumption $G = A.F$ means $G(x) = F \big( A^{-1}(x) \big)$ for all $x \in \bR$.
  Therefore $G \big( A(x_1) \big) = F(x_1) < F(x_2) = G \big( A(x_2) \big)$.
  We also get $A(x_j) = \inf \big\{ y \in \bR \; \big| \; G(y) = F(x_j) \big\}$ for
  $j=1,2$ by strict monotonicity and bijectivity of $A$
  (if, for example, there would exist a $y'_2 < A(x_2)$ such that $G(y'_2) = F(x_2)$,
  then the point $x'_2 = A^{-1}(y'_2) < x_2$ would be such that
  $F\big(A^{-1}(y'_2)\big) = G(y'_2) = F(x_2)$, contradicting the minimality of $x_2$).

  If $\widetilde{A} \in \OriAffR$ is also such that $G = \widetilde{A}.F$, then
  the same holds for it:
  $\widetilde{A}(x_j) = \inf \big\{ y \in \bR \; \big| \; G(y) = F(x_j) \big\}$
  for $j=1,2$. We conclude that
  \begin{align*}
    \widetilde{A}(x_1) \, = \; &
        \inf \big\{ y \in \bR \; \big| \; G(y) = F(x_1) \big\} \; = \, A(x_1) \\
    \widetilde{A}(x_2) \, = \; &
        \inf \big\{ y \in \bR \; \big| \; G(y) = F(x_2) \big\} \; = \, A(x_2) .
  \end{align*}
  But an affine map of $\bR$ is determined by its values at two distinct points:
  from $a x_1 + b = y_1$ and $a x_2 + b = y_2$ with $x_1 \ne x_2$ one can solve $a, b$.
  Therefore we must have $\widetilde{A} = A$.
\end{proof}

\begin{lemma}[Degeneration by shrinking affine transformations]
  \label{lem:degenerate-shrinking-limit}
  \uses{def:cdf-degenerate, lem:oriented-affine-action-on-cdf}
  \lean{CumulativeDistributionFunction.isDegenerate_of_tendsto_shrinking}
  \leanok
  Let $(F_n)_{n \in \bN}$ be a sequence of c.d.f.s which converges to a c.d.f. $G$,
  $F_n \tendsinlaw G$. Consider affine transformations of the form $A_n(x) = a_n x + b_n$,
  with $a_n > 0$ and $b_n \in \bR$, such that $a_n \to 0$ and $b_n \to \beta \in \bR$
  as $n \to \infty$. Then $A_n . F_n \tendsinlaw \widetilde{G}$, where
  $\widetilde{G}$ is the degenerate c.d.f. of the delta mass at $\beta$.
\end{lemma}
\begin{proof}
  % \uses{}
  % \leanok
  It suffices to prove that for any $x < \beta$ we have $\widetilde{G}(x) = 0$
  and for any $x > \beta$ we have $\widetilde{G}(x) = 1$.

  Let us focus on the latter, so let $x > \beta$. Let also $\eps > 0$; we will prove
  that $\widetilde{G}(x) > 1 - \eps$, and the claim will follow.

  By density of continuity points of $\widetilde{G}$, we can choose $x'$ such that
  $\beta < x' < x$ and $\widetilde{G}$ is continuous at $x'$. Then the assumed
  convergence $A_n . F_n \tendsinlaw \widetilde{G}$ implies that
  $(A_n . F_n)(x') \to \widetilde{G}(x')$. Since $\widetilde{G}(x) \ge \widetilde{G}(x')$,
  it suffices to prove that $\widetilde{G}(x') > 1 - \eps$.

  Since $G$ is a c.d.f., we can choose a continuity point $z$ of $G$
  large enough so that $G(z) > 1 - \eps$.
  Then by the assumed convergence $F_n \tendsinlaw G$, we have
  $F_n(z) \to G(z)$. By definition we have
  \begin{align*}
    (A_n . F_n)(A_n(z)) \; = \; F_n \big( A_n^{-1} (A_n(z)) \big) \; = \; F_n(z)
      \; \longrightarrow \; G(z) . %\; > \;  1-\eps .
  \end{align*}
  Note that $A_n(z) = a_n z + b_n \to \beta$ as $n \to \infty$ by the assumptions
  $a_n \to 0$ and $b_n \to \beta$. In particular, for $n$ large enough,
  we have $A_n(z) < x'$.
  Therefore, for $n$ large enough
  \begin{align*}
    F_n(z) \; = \; (A_n . F_n) \big( A_n(z) \big) \; \le \; (A_n . F_n)(x') .
  \end{align*}
  The LHS tends to $G(z)$ and the RHS tends to $\widetilde{G}(x')$,
  showing
  \begin{align*}
    1 - \eps \; < \; G(z) \; \le \; \widetilde{G}(x') \; \le \; \widetilde{G}(x) .
  \end{align*}
  This concludes the proof that $\widetilde{G}(x) = 1$ for all $x > \beta$.

  The proof that $\widetilde{G}(x) = 0$ for all $x < \beta$ is similar.
\end{proof}

\begin{lemma}[Impossibility of expanding affine transformations]
  \label{lem:impossible-expanding-limit}
  \uses{def:cdf, lem:oriented-affine-action-on-cdf}
  \lean{CumulativeDistributionFunction.not_tendsto_cdf_of_expanding_of_tendsto_not_isDegenerate}
  \leanok
  Let $(F_n)_{n \in \bN}$ be a sequence of c.d.f.s which converges to a nondegenerate
  c.d.f. $G$, $F_n \tendsinlaw G$. Consider affine transformations of the form $A_n(x) = a_n x + b_n$,
  with $a_n > 0$ and $b_n \in \bR$, such that $a_n \to +\infty$
  as $n \to \infty$. Then $A_n . F_n$ cannot converge to any c.d.f.
\end{lemma}
\begin{proof}
  \uses{lem:exists-two-nontrivial-continuity-pts-cdf}
  \leanok
  Since $G$ is assumed nondegenerate, by Lemma~\ref{lem:exists-two-nontrivial-continuity-pts-cdf}
  we can pick two continuity points $x_1 < x_2$
  of $G$ such that $0 < G(x_1) \le G(x_2) < 1$. Then from the assumption
  $F_n \tendsinlaw G$ we get $F_n(x_1) \to G(x_1)$ and $F_n(x_2) \to G(x_2)$.

  Assume, by way of contradiction, that we have convergence
  $A_n . F \tendsinlaw \widetilde{G}$ to some c.d.f. $\widetilde{G}$.

  We claim that then $\big(A_n(x_1)\big)_{n \in \bN}$ is bounded from below and
  $\big(A_n(x_2)\big)_{n \in \bN}$ is bounded from above.
  Since
  \begin{align*}
    a_n = \frac{A_n(x_2) - A_n(x_1)}{x_2 - x_1} ,
  \end{align*}
  this will show that $(a_n)_{n \in \bN}$ is bounded from above, contradicting
  the assumption $a_n \to +\infty$, and finishing the proof.

  To show that $\big(A_n(x_1)\big)_{n \in \bN}$ is bounded from below,
  choose a continuity point $z$ of $\widetilde{G}$ such that
  $\widetilde{G}(z) < G(x_1)$. Then the assumed convergence
  $A_n . F \tendsinlaw \widetilde{G}$ implies
  $(A_n . F)(z) \to \widetilde{G}(z)$.
  On the other hand, if $\big(A_n(x_1)\big)_{n \in \bN}$ is not
  bounded from below, then along some subsequence $(n_k)_{k \in \bN}$ of indices
  we have $A_{n_k}(x_1) < z$, and for those indices we then have
  \begin{align*}
    F_{n_k}(x_1) \; = \; (A_{n_k} . F_{n_k}) \big( A_{n_k}(x_1) \big)
      \; \le \; (A_{n_k} . F_{n_k})(z) .
  \end{align*}
  The LHS tends to $G(x_1)$ as $k \to \infty$, whereas the RHS tends
  to $\widetilde{G}(z)$. We get $G(x_1) \le \widetilde{G}(z)$, contradicting
  the choice of $z$. This shows that $\big(A_n(x_1)\big)_{n \in \bN}$ must
  in fact be bounded from below.

  The proof that $\big(A_n(x_2)\big)_{n \in \bN}$ must be bounded from above
  is similar.
\end{proof}

\begin{theorem}[Convergence to types]
  \label{thm:convergence-to-types}
  \uses{def:degenerate-cdf, lem:oriented-affine-action-on-cdf, def:cdf-type}
  % \lean{}
  % \leanok
  Suppose that $(F_n)_{n \in \bN}$ is a sequence of
  c.d.f.s which converges to a nondegenerate c.d.f. $G$, i.e.,
  $F_n \tendsinlaw G$ as $n \to \infty$.
  Let $(A_n)_{n \in \bN}$ be a sequence of oriented
  affine isomorphisms of $\bR$, $A_n \in \OriAffR$
  such that $A_n.F_n \tendsinlaw \widetilde{G}$ for some c.d.f.
  $\widetilde{G}$.

  If we write $A_n(x) = a_n x + b_n$, then $(a_n)_{n \in \bN}$
  and $(b_n)_{n \in \bN}$ are bounded sequences.

  If $\widetilde{G}$ is nondegenerate, then
  $A_n \to A \in \OriAffR$ and $A.G = \widetilde{G}$.
  In particular $G$ and $\widetilde{G}$ are of the same type.
  Moreover, $A$ is the unique affine transformation for which
  the equality $A.G = \widetilde{G}$ holds.
\end{theorem}
\begin{proof}
  \uses{lem:impossible-expanding-limit, lem:degenerate-shrinking-limit,
    lem:unique-affine-relation-cdf, lem:action-on-cdf-continuous}
  % \leanok
  Let us first argue that $(a_n)_{n \in \bN}$ are bounded.
  If not, then by passing to a subsequence, we have $a_{n_k} \to +\infty$.
  But since $F_{n_k} \tendsinlaw G$ and $G$ is nondegenerate,
  it contradicts Lemma~\ref{lem:impossible-expanding-limit}
  to have $A_{n_k} . F_{n_k} \tendsinlaw \widetilde{G}$.
  Therefore $(a_n)_{n \in \bN}$ must be bounded: there exists
  some $M>0$ such that $a_n \le M$ for all $n \in \bN$.

  Let us then argue that $(b_n)_{n \in \bN}$ is bounded.
  If not, then we can extract a subsequence such
  that either $b_{n_k} \to -\infty$ or $b_{n_k} \to +\infty$.
  Let us prove the impossibility of the second one, the first is similar.
  So assume that $b_{n_k} \to +\infty$.
  Since $G$ is nondegenerate, we may pick a continuity point $x_0$
  of $G$ such that $0 < G(x_0) < 1$. Then we have $F_n(x_0) \to G(x_0)$
  by the assumption $F_n \tendsinlaw G$.
  We may also pick a continuity point $z$ of $\widetilde{G}$
  such that $\widetilde{G}(z) > G(x_0)$. Then we have
  $(A_n.F_n)(z) \to \widetilde{G}(z)$
  by the assumption $A_n.F_n \tendsinlaw \widetilde{G}$.
  But $A_{n_k}(x_0) = a_{n_k} x_0 + b_{n_k} \to +\infty$, since $0 < a_{n_k} \le M$
  and $b_{n_k} \to +\infty$. Therefore we have for all large enough $k$
  that $A_{n_k}(x_0) > z$. And then
  \begin{align*}
    (A_{n_k}.F_{n_k})(z) \; \le \; (A_{n_k}.F_{n_k})\big(A_{n_k}(x_0)\big) \; = \; F_{n_k}(x_0) .
  \end{align*}
  The LHS tends to $\widetilde{G}(z)$ as $k \to \infty$, and the RHS tends to $G(x_0)$.
  Therefore we get $\widetilde{G}(z) \le G(x_0)$, contradicting the choice of $z$.
  This shows that we cannot have $b_{n_k} \to +\infty$. Similarly one proves that
  we cannot have $b_{n_k} \to -\infty$. We conclude that $(b_n)_{n \in \bN}$
  is indeed bounded.

  From now on, suppose furthermore that also $\widetilde{G}$ is nondegenerate.
  We claim that then $(a_n)_{n \in \bN}$ is bounded away from $0$:
  for some $\eps > 0$ we have $a_n \ge \eps$ for all $n \in \bN$.
  If not, then we could extract a subsequence such that $a_{n_k} \to 0$
  and also $b_n \to \beta$ (since $b_n$ is bounded).
  But since $F_{n_k} \tendsinlaw G$ and $G$ is nondegenerate,
  from Lemma~\ref{lem:degenerate-shrinking-limit}
  we would get $A_{n_k} . F_{n_k} \tendsinlaw \widetilde{G}_0$,
  where $\widetilde{G}_0$ is degenerate. But by assumption
  $A_{n_k} . F_{n_k} \tendsinlaw \widetilde{G}$, where
  $\widetilde{G}$ is nondegenerate; this is impossible by uniqueness
  of limits for convergence in distribution. Therefore
  $(a_n)_{n \in \bN}$ must indeed be bounded away from $0$.

  Note that since $(a_n)_{n \in \bN}$ is bounded away from $0$ and $+\infty$,
  and $(b_n)_{n \in \bN}$ is bounded, we can extract a subsequence such that
  $a_{n_k} \to \alpha \in (0,+\infty)$ and $b_{n_k} \to \beta \in \bR$.
  By assumption, we have $A_{n_k} . F_{n_k} \tendsinlaw \widetilde{G}$.
  But since $A_{n_k} \to A$ where $A(x) = \alpha x + \beta$
  and we have also assumed $F_{n_k} \tendsinlaw G$, this implies by
  continuity (Lemma~\ref{lem:action-on-cdf-continuous}) that $A_{n_k} . F_{n_k} \tendsinlaw A . G$.
  By uniqueness of limits, we get $A . G = \widetilde{G}$.

  To prove that in fact $A_n \to A$, not just along a subsequence, note the
  following. From any subsequence $A_{n_k}$, we can extract a further
  convergent subsequence of values of $a_{n_{k_\ell}}$ and $b_{n_{k_\ell}}$
  values as above, with limits $\alpha' \in (0,+\infty)$
  and $\beta' \in \bR$. The same argument as above shows that
  $A' . G = \widetilde{G}$ where $A'(x) = \alpha' x + \beta'$.
  Lemma~\ref{lem:unique-affine-relation-cdf} then says that we must have
  $A' = A$, i.e., $\alpha' = \alpha$ and $\beta' = \beta$. Since any
  subsequence has a convergent further subsequence with the same limit,
  the entire sequence must converge, $A_n \to A$.
  The proof is complete.
\end{proof}

\begin{theorem}[Convergence to types again]
  \label{thm:convergence-to-types-different-affine}
  \uses{def:cdf, lem:oriented-affine-action-on-cdf, def:cdf-type}
  % \lean{} % also `AffineIncrEquiv.inv_mul_eq_mkOfCoefs`
  % \leanok
  Let $(A_n)_{n \in \bN}$ and
  $(\widetilde{A}_n)_{n \in \bN}$ be two sequences of oriented
  affine isomorphisms of $\bR$, $A_n, \widetilde{A}_n \in \OriAffR$.
  Write $A_n(x) = a_n x + b_n$
  and $\widetilde{A}_n(x) = \tilde{a}_n x + \tilde{b}_n$,
  and for the inverses $A_n^{-1}(x) = c_n x + d_n$
  and $\widetilde{A}^{-1}_n(x) = \tilde{c}_n x + \tilde{d}_n$.

  Let $(F_n)_{n \in \bN}$ be a sequence of c.d.f.s
  such that $A_n.F_n \tendsinlaw G$, with
  $G$ a nondegenerate c.d.f.

  Then the convergence of also $\widetilde{A}_n.F_n \tendsinlaw G$
  holds if and only if the coefficients of the affine maps
  satisfy the relations
  \begin{align*}
    \frac{\tilde{a}_n}{a_n} \to 1
    \quad \text{ and } \quad
    \frac{a_n \tilde{b}_n - \tilde{a}_n b_n}{a_n} \to 0,
  \end{align*}
  or equivalently,
  \begin{align*}
    \frac{\tilde{c}_n}{c_n} \to 1
    \quad \text{ and } \quad
    \frac{\tilde{d}_n - d_n}{c_n} \to 0 .
  \end{align*}
\end{theorem}
\begin{proof}
  \uses{thm:convergence-to-types}
  % \leanok
  \emph{(This is actually just a special case of what is stated as
  Corollary~\ref{cor:convergence-to-types-different-limit} below.
  The better organization is to prove that corollary
  directly, and obtain this lemma as a special case.)}

  We will apply the convergence to types with the reference sequence
  $(A_n.F_n)_{n \in \bN}$, which by assumption tends to a
  nondegenerate $G$.

  To express the other sequence in terms of the reference sequence, we write
  \begin{align*}
    \widetilde{A}_n.F_n \; = \; (\widetilde{A}_n A_n^{-1}).(A_n.F_n) .
  \end{align*}
  By assumption this also tends to $G$.

  Theorem~\ref{thm:convergence-to-types} applies, and guarantees convergence
  $\widetilde{A}_n A_n^{-1} \to A$ to some $A \in \OriAffR$,
  and it also implies $A.G = G$ (note that the other limit is also $G$
  by our assumptions).
  However, the unique $A$ for which we have $A.G = G$ is $A = \mathrm{id}$.
  We therefore get $\widetilde{A}_n A_n^{-1} \to \id$.
  To explicitly see the coefficients, note that
  \begin{align*}
    A_n^{-1}(x) \; = \; & a_n^{-1} (x - b_n)
  \end{align*}
  and
  \begin{align*}
    \widetilde{A}_n \big( A_n^{-1}(x) \big)
    \; = \; & \tilde{a}_n \big(a_n^{-1} (x - b_n)\big) + \tilde{b}_n \\
    \; = \; & \tilde{a}_n a_n^{-1} x - \tilde{a}_n a_n^{-1} b_n + \tilde{b}_n .
  \end{align*}
  The convergence $\widetilde{A}_n A_n^{-1} \to \id$ is equivalent to
  the convergence of the coefficients,
  \begin{align*}
    \tilde{a}_n a_n^{-1} \; \longrightarrow \; & 1 \\
    - \tilde{a}_n a_n^{-1} b_n + \tilde{b}_n \; \longrightarrow \; & 0 .
  \end{align*}
  The second one can be rewritten as
  \begin{align*}
    \frac{a_n \tilde{b}_n - \tilde{a}_n b_n}{a_n} \; \longrightarrow \; & 0 .
  \end{align*}
\end{proof}

\begin{corollary}[Convergence to types with different limits]
  \label{cor:convergence-to-types-different-limit}
  \uses{def:cdf, lem:oriented-affine-action-on-cdf, def:cdf-type}
  % \lean{} % also `AffineIncrEquiv.inv_mul_eq_mkOfCoefs`
  % \leanok
  Let $(A_n)_{n \in \bN}$ and
  $(\widetilde{A}_n)_{n \in \bN}$ be two sequences of oriented
  affine isomorphisms of $\bR$, $A_n, \widetilde{A}_n \in \OriAffR$.
  Write $A_n(x) = a_n x + b_n$
  and $\widetilde{A}_n(x) = \tilde{a}_n x + \tilde{b}_n$,
  and for the inverses $A_n^{-1}(x) = c_n x + d_n$
  and $\widetilde{A}^{-1}_n(x) = \tilde{c}_n x + \tilde{d}_n$.

  Let $(F_n)_{n \in \bN}$ be a sequence of c.d.f.s
  such that $A_n.F_n \tendsinlaw G$ and
  $\widetilde{A}_n.F_n \tendsinlaw \widetilde{G}$, with
  $G$ and $\widetilde{G}$ nondegenerate c.d.f.s.
  Then for some $\alpha>0$ and $\beta \in \bR$ we have
  \begin{align*}
    \frac{\tilde{a}_n}{a_n} \to \alpha
    \quad \text{ and } \quad
    \frac{a_n \tilde{b}_n - \tilde{a}_n b_n}{a_n} \to \beta,
  \end{align*}
  and we have
  \begin{align*}
    A.G = \widetilde{G}
    \qquad \text{ where } \; A(x) = \alpha x + \beta .
  \end{align*}
  Equivalently, with $\gamma = \alpha^{-1}$ and $\delta = -\alpha^{-1} \beta$
  so that $A^{-1}(x) = \gamma x + \delta$, we have
  \begin{align*}
    \frac{\tilde{c}_n}{c_n} \to \gamma
    \quad \text{ and } \quad
    \frac{\tilde{d}_n - d_n}{c_n} \to \delta .
  \end{align*}

  In particular, $G$ and $\widetilde{G}$ have the same type.
\end{corollary}
\begin{proof}
  \uses{thm:convergence-to-types-different-affine}
  % \leanok
  We will apply the convergence to types with the reference sequence
  $(A_n.F_n)_{n \in \bN}$, which by assumption tends to a
  nondegenerate $G$.

  To express the other sequence in terms of the reference sequence, we write
  \begin{align*}
    \widetilde{A}_n.F_n \; = \; (\widetilde{A}_n A_n^{-1}).(A_n.F_n) .
  \end{align*}
  By assumption this tends to a nondegenerate $\widetilde{G}$.

  Theorem~\ref{thm:convergence-to-types} applies, and guarantees convergence
  $\widetilde{A}_n A_n^{-1} \to A$ to some $A \in \OriAffR$,
  and it also implies $A.G = \widetilde{G}$.

  Write $A(x) = \alpha x + \beta$. To explicitly see the coefficients of
  $\widetilde{A}_n A_n^{-1}$, note that
  \begin{align*}
    A_n^{-1}(x) \; = \; & a_n^{-1} (x - b_n)
  \end{align*}
  and
  \begin{align*}
    \widetilde{A}_n \big( A_n^{-1}(x) \big)
    \; = \; & \tilde{a}_n \big(a_n^{-1} (x - b_n)\big) + \tilde{b}_n \\
    \; = \; & \tilde{a}_n a_n^{-1} x - \tilde{a}_n a_n^{-1} b_n + \tilde{b}_n .
  \end{align*}
  The convergence $\widetilde{A}_n A_n^{-1} \to A$ is equivalent to
  the convergence of the coefficients,
  \begin{align*}
    \tilde{a}_n a_n^{-1} \; \longrightarrow \; & \alpha \\
    - \tilde{a}_n a_n^{-1} b_n + \tilde{b}_n \; \longrightarrow \; & \beta .
  \end{align*}
  The second one can be rewritten as
  \begin{align*}
    \frac{a_n \tilde{b}_n - \tilde{a}_n b_n}{a_n} \; \longrightarrow \; & \beta .
  \end{align*}
\end{proof}

\begin{lemma}[A choice of normalizing constants for convergence to types]
  \label{lemma:convergence-to-types-normalization}
  \uses{def:cdf, lem:oriented-affine-action-on-cdf, def:cdf-type, def:lc-inverse}
  % \lean{}
  % \leanok
  (It is possible to choose normalization constants for the affine
  transformations using the left-continuous inverses of the c.d.f.s.
  TODO: Precise statement.)
%  Suppose that $(F_n)_{n \in \bN}$ is a sequence of c.d.f.s
%  and $(A_n)_{n \in \bN}$ %and $(\widetilde{A}_n)_{n \in \bN}$
%  be a sequence of oriented
%  affine isomorphisms of $\bR$, $A_n \in \OriAffR$,
%  such that $A_n.F_n \tendsinlaw G$ where
%  $G$ is a nondegenerate c.d.f.
%
%  (TODO: The following is not quite correct. Need to assume one chooses
%  continuity points of the left-continuous inverse.)
%
%  Then for any $0 < p_1 < p_2 < 1$ and $0 < p_0 < 1$, defining
%  \begin{align*}
%    \tilde{a}_n = \frac{1}{\lcinv{F}_n(p_2) - \lcinv{F}_n(p_1)}
%    \quad \text{ and } \quad
%    \tilde{b}_n = \frac{\lcinv{F}_n(p_0)}{\lcinv{F}_n(p_2) - \lcinv{F}_n(p_1)}
%  \end{align*}
%  we have that
%  \begin{align*}
%    \widetilde{A}_n.F_n \tendsinlaw \widetilde{G}
%    \qquad \text{ where } \; \widetilde{A}_n(x) = \tilde{a}_n x + \tilde{b}_n .
%  \end{align*}
\end{lemma}
\begin{proof}
  \uses{thm:convergence-to-types-different-affine}
  % \leanok
  \ldots
\end{proof}


\section{One-parameter subgroups of affine transformations}

\begin{definition}[Subgroup of translations]
  \label{def:translation-subgroup}
  \uses{def:oriented-affine-isomorphism}
  \lean{AffineIncrEquiv.homOfIndex₀, AffineIncrEquiv.subGroupOfIndex₀}
  % \leanok
  The mapping $s \mapsto A_s$ with
  \begin{align*}
    A_s(x) = x + s
  \end{align*}
  is a homomorphism $\bR \to \OriAffR$. The image of this
  homomorphism is the \emph{subgroup of translations} in $\OriAffR$.
\end{definition}

\begin{lemma}[Only translations have no fixed points]
  \label{lem:no-fixed-point-implies-translation}
  \uses{def:translation-subgroup}
  \lean{AffineIncrEquiv.mem_subGroupOfIndex₀_of_no_fixed_point}
  \leanok
  If $A \in \OriAffR$ has no fixed points (no $x \in \bR$ such that $A(x) = x$)
  then $A$ belongs to the subgroup of translations, i.e., $A(x) = x + s$
  for some $s \in \bR$ (in fact $s \ne 0$).
\end{lemma}
\begin{proof}
  % \uses{}
  % \leanok
  Let us prove this by contrapositive: that any element $A$ which is not a
  translation must have a fixed point.
  So assume that $A$ is not a translation, i.e.,
  $A(x) = a x + b$ with some $a \ne 1$. Then the fixed point equation
  $A(x) = x$ reads
  \begin{align*}
    a x + b = x  ,
  \end{align*}
  and it has a solution $x = \frac{-b}{a-1} \in \bR$, which then is a
  fixed point of $A$.
\end{proof}

\begin{lemma}[Conjugate of translation is translation]
  \label{lem:conjugate-translation}
  \uses{def:translation-subgroup}
  \lean{AffineIncrEquiv.conjugate_homOfIndex₀}
  % \leanok
  Let $A^{(\beta)}_s = x + \beta s$ for $s, \beta \in \bR$ as in
  Definition~\ref{def:translation-subgroup}. Let
  also $B \in \OriAffR$ be given by $B(x) = a x + b$. Then
  \begin{align*}
    B \, A^{(\beta)}_s \, B^{-1} = A^{(a \beta)}_{s} .
  \end{align*}
\end{lemma}
\begin{proof}
  % \uses{}
  % \leanok
  Calculate, for $x \in \bR$
  \begin{align*}
    (B A^{(\beta)}_s B^{-1})(x) \; = \; & (B A^{(\beta)}_s)\big( \frac{x-b}{a} \big) \\
    \; = \; & B\big( \frac{x-b}{a} + \beta s \big) \\
    \; = \; & a \big( \frac{x-b}{a} + \beta s \big) + b \\
    \; = \; & x - b + a \beta s + b \\
    \; = \; & x + a \beta s \; = \; A^{(a \beta)}_s(x).
  \end{align*}
\end{proof}

\begin{definition}[Subgroup fixing a point]
  \label{def:fixing-subgroup}
  \uses{def:oriented-affine-isomorphism}
  \lean{AffineIncrEquiv.homOfIndex, AffineIncrEquiv.subGroupOfIndex}
  % \leanok
  The mapping $s \mapsto A_s$ with
  \begin{align*}
    A_s(x) = e^{s} (x - c) + c
  \end{align*}
  is a homomorphism $\bR \to \OriAffR$. The image of this
  homomorphism is the \emph{subgroup fixing $c$} in $\OriAffR$.
\end{definition}

\begin{lemma}[Characterization of the subgroup fixing a point]
  \label{lem:fixing-subgroup-characterization}
  \uses{def:fixing-subgroup}
  \lean{AffineIncrEquiv.mem_subGroupOfIndex_iff_fixed_point}
  \leanok
  An orientation-preserving affine transformation $A \in \OriAffR$
  belongs to the subgroup fixing $c \in \bR$ if and only if $A(c) = c$.

  \emph{(Note that the subgroup is a priori defined as the image of a homomorphism,
  so the statement indeed requires a proof.)}
\end{lemma}
\begin{proof}
  % \uses{}
  % \leanok
  Suppose first that $A$ is an element of the said subgroup, i.e.,
  $A(x) = e^{s} (x - c) + c$ for some $s \in \bR$. Then clearly $A(c) = c$.

  Suppose then that $A(c) = c$. Write $A(x) = a x + b$ for $a>0$ and $b \in \bR$.
  Plug in $x=c$ in the assumed fixed point property to obtain
  \begin{align*}
    a c + b = c .
  \end{align*}
  The above can be solved to give $b = (1-a)c$.
  Also since $a>0$, we can write $a = e^s$ with $s \in \bR$.
  With these, the formula for $A$ simplifies to
  \begin{align*}
    A(x) = e^s x + c (1 - e^s) = e^s \big( x - c \big) + c .
  \end{align*}
  This shows $A = A_s$ as desired (with $A_s$ as in Definition~\ref{def:fixing-subgroup}).
\end{proof}

\begin{lemma}[Conjugate of fixing is fixing image]
  \label{lem:conjugate-fixing}
  \uses{def:fixing-subgroup}
  \lean{AffineIncrEquiv.conjugate_homOfIndex}
  % \leanok
  Let $A^{(\alpha;c)}_s = e^{\alpha s} (x - c) + c$ for $\alpha, c \in \bR$ as in
  Definition~\ref{def:fixing-subgroup}. Let
  also $B \in \OriAffR$ be given by $B(x) = a x + b$. Then
  \begin{align*}
    B \, A^{(\alpha;c)}_s \, B^{-1} = A^{(\alpha;B(c))}_{s} .
  \end{align*}
\end{lemma}
\begin{proof}
  % \uses{}
  % \leanok
  Calculate, for $x \in \bR$
  \begin{align*}
    (B A^{(\alpha;c)}_s B^{-1})(x) \; = \; & (B A^{(\alpha;c)}_s)\big( \frac{x-b}{a} \big) \\
    \; = \; & B\big( e^{\alpha s} (\frac{x-b}{a} - c) + c \big) \\
    \; = \; & a e^{\alpha s} (\frac{x-b}{a} - c) + a c + b \\
    \; = \; & e^{\alpha s} \big( x - (a c + b) \big) + (a c + b) \\
    \; = \; & e^{\alpha s} \big( x - B(c) \big) + B(c) \\
    \; = \; & A^{(\alpha;B(c))}_s (x) .
  \end{align*}
\end{proof}


\section{Self-similarity characterizations of the extreme value distributions}

\begin{lemma}[Continuous parameter extreme value limit relation]
  \label{lem:continuous-parameter-ev-limit}
  \uses{def:cdf, lem:oriented-affine-action-on-cdf}
  \lean{continuous_parameter_ev_limit_relation}
  \leanok
  Let $F$ be a c.d.f.

  \emph{(Note that below we use
  the sequence $(F^n)_{n \in \bN}$ of $n$th powers of a fixed c.d.f.,
  not a sequence of arbitrary c.d.f.s.
  Recall that the $n$th power $F^n$ is the c.d.f. of the maximum
  of $n$ independent random variables with the distribution $F$.)}

  Suppose that for a sequence $(A_n)_{n \in \bN}$ of oriented
  affine isomorphisms of $\bR$, $A_n \in \OriAffR$, we have
  \begin{align*}
    A_n.F^n \tendsinlaw G ,
  \end{align*}
  where $G$ is a c.d.f.

  Then, for any $t > 0$, denoting by $G^t$ the c.d.f. given by
  $G^t(x) = \big( G(x) \big)^t$, we have
  \begin{align*}
    A_{n}.F^{\lfloor n t \rfloor} \tendsinlaw G^t ,
  \end{align*}
  where, for $x \in \bR$, the floor notation $\lfloor x \rfloor$ stands
  for the greatest integer $k \in \bZ$ such that $k \le x$.
\end{lemma}
\begin{proof}
  % \uses{}
  % \leanok
  Let $t > 0$ and let $x \in \bR$ be a continuity point of $G$.
  For $n \in \bN$, calculate
  \begin{align*}
      \big((A_n . F) (x) \big)^{\lfloor n t \rfloor}
    %\; = \; \big(F(A_n^{-1}(x)) \big)^{\lfloor n t \rfloor} \\
    \; = \; & \Big( \big((A_n . F) (x) \big)^n \Big)^{\lfloor n t \rfloor / n} .
  \end{align*}
  By assumption, we have $\big((A_n . F) (x) \big)^n \to G(x)$ as $n \to \infty$.
  Also $\lfloor n t \rfloor / n \to t$ as $n \to \infty$.
  By (joint) continuity of the power function
  $(x,y) \mapsto x^y = \exp\big( y \, \log(x) \big)$,
  we get that the expression above tends to $G(x)^t$.

  Finally noting that the continuity points of $G^t$ are the same as the
  continuity points of $G$, the above in fact proves the asserted
  $A_{n}.F^{\lfloor n t \rfloor} \tendsinlaw G^t$.
\end{proof}

\begin{lemma}[Self-similarity of extreme value distributions]
  \label{lem:self-similarity-of-extreme-value-distributions}
  \uses{def:cdf, lem:oriented-affine-action-on-cdf, def:extr-val-distr}
  \lean{CumulativeDistributionFunction.IsExtremeValueDistr.self_similar}
  \leanok
  %Let $F$ be a c.d.f.
  %
  %\emph{(Note that below we use
  %the sequence $(F^n)_{n \in \bN}$ of $n$th powers of a fixed c.d.f.,
  %not a sequence of arbitrary c.d.f.s.
  %Recall that the $n$th power $F^n$ is the c.d.f. of the maximum
  %of $n$ independent random variables with the distribution $F$.)}
  %
  %Suppose that for a sequence $(A_n)_{n \in \bN}$ of oriented
  %affine isomorphisms of $\bR$, $A_n \in \OriAffR$, we have
  %\begin{align*}
  %  A_n.F^n \tendsinlaw G ,
  %\end{align*}
  %where $G$ is a c.d.f.
  Suppose that $G$ is an extreme-value distribution.
  Then there exists a family $(A_t)_{t > 0}$ of
  oriented affine isomorphisms of $\bR$, $A_t \in \OriAffR$,
  such that for any $t > 0$
  \begin{align*}
    G^t = A_t . G .
  \end{align*}

  Moreover, $t \mapsto A_t$ is a measurable homomorphism
  of multiplicative groups $(0,+\infty) \to \OriAffR$.
\end{lemma}
\begin{proof}
  \uses{lem:continuous-parameter-ev-limit, thm:convergence-to-types}
  % \leanok
  \ldots
\end{proof}

%\begin{lemma}[Continuous parameter limit relation]
%  \label{lem:continuous-parameter-ev-limit}
%  \uses{def:cdf, lem:oriented-affine-action-on-cdf}
%  % \lean{}
%  % \leanok
%  Let $F$ be a c.d.f.
%
%  Suppose that for a sequence $(A_n)_{n \in \bN}$ of oriented
%  affine isomorphisms of $\bR$, and for a family $(G_t)_{t > 0}$
%  of c.d.f.s, \ldots.
%
%
%  and any $t > 0$, we have
%  we have
%  \begin{align*}
%    A_n.F^n \tendsinlaw G_t ,
%  \end{align*}
%  where $G_t$ is a nondegenerate c.d.f.
%
%  Then, for any $t > 0$, denoting by $G^t$ the c.d.f. given by
%  $G^t(x) = \big( G(x) \big)^t$, we have
%  \begin{align*}
%    A_{n}.F^{\lfloor n t \rfloor} \tendsinlaw G^t .
%  \end{align*}
%\end{lemma}

%%\begin{lemma}[Functional equation in one parameter subgroups of affine isomorphisms]
%%  \label{lem:affine-one-parameter-subgroup}
%%  \uses{def:lem:oriented-affine-action-on-cdf}
%%  % \lean{}
%%  % \leanok
%%  Suppose that $(A_t)_{t > 0}$ is a family of affine isomorphisms of $\bR$
%%  such that
%%  \begin{align*}
%%    A_{s t} = A_s \circ A_t  \qquad \text{ for any } s, t > 0 .
%%  \end{align*}
%%  Write $A_t(x) = a_t x + b_t$, with $a_t > 0$ and $b_t \in \bR$.
%%  Then we have, for any $s, t > 0$,
%%  \begin{align*}
%%    a_{t s} = \; & a_t \, a_s
%%      \qquad \text{ and } \\
%%    b_{t s} = \; & a_t \, b_s + b_t .
%%  \end{align*}
%%  (Also by symmetry $b_{t s} = a_s \, b_t + b_s$.)
%%\end{lemma}
%%\begin{proof}
%%  % \uses{}
%%  % \leanok
%%\end{proof}
%%
%%\begin{lemma}[Functional equation scaling coefficient solution]
%%  \label{lem:solution-functional-eqn-scaling}
%%  %\uses{def:lem:oriented-affine-action-on-cdf}
%%  % \lean{}
%%  % \leanok
%%  Suppose that $a \colon (0,+\infty) \to (0,+\infty)$ is a
%%  measurable function satisfying, for any $s, t > 0$,
%%  \begin{align*}
%%    a(t s) = \; & a(t) \, a(s) .
%%  \end{align*}
%%  Then there exists a $\theta \in \bR$ such that for all $t > 0$,
%%  \begin{align*}
%%    a(t) = t^{\theta} .
%%  \end{align*}
%%\end{lemma}
%%\begin{proof}
%%  % \uses{}
%%  % \leanok
%%\end{proof}
%%
%%\begin{lemma}[Functional equation translation coefficient solution with $\theta = 0$]
%%  \label{lem:solution-functional-eqn-translation-no-scaling}
%%  %\uses{def:lem:oriented-affine-action-on-cdf}
%%  % \lean{}
%%  % \leanok
%%  Suppose that
%%  $b \colon (0,+\infty) \to \bR$ is a measurable function
%%  satisfying, for any $s, t > 0$,
%%  \begin{align*}
%%    b(t s) = \; & b(s) + b(t) .
%%  \end{align*}
%%  Then there exists a constant $c$ such that for $t \in (0,+\infty)$
%%  we have
%%  \begin{align*}
%%    b(t) = \; & -c \, \log (t) .
%%  \end{align*}
%%\end{lemma}
%%\begin{proof}
%%  % \uses{}
%%  % \leanok
%%\end{proof}
%%
%%\begin{lemma}[Functional equation translation coefficient solution with $\theta \ne 0$]
%%  \label{lem:solution-functional-eqn-translation-with-scaling}
%%  %\uses{def:lem:oriented-affine-action-on-cdf}
%%  % \lean{}
%%  % \leanok
%%  Suppose that $\theta \in \bR \setminus 0$ and
%%  $b \colon (0,+\infty) \to \bR$ is a measurable function
%%  satisfying, for any $s, t > 0$,
%%  \begin{align*}
%%    b(t s) = \; & t^{\theta} \, b(s) + b(t) .
%%  \end{align*}
%%  Then there exists a constant $c$ such that for $t \in (0,+\infty) \setminus \set{1}$
%%  we have
%%  \begin{align*}
%%    b(t) = \; & c (1 - t^{-\theta}) .
%%  \end{align*}
%%\end{lemma}
%%\begin{proof}
%%  % \uses{}
%%  % \leanok
%%\end{proof}

\begin{lemma}[Self-similar continuous c.d.f. family characterization $\gamma = 0$]
  \label{lem:characterization-self-similar-family-zero-index}
  \uses{def:lem:oriented-affine-action-on-cdf, def:translation-subgroup}
  \lean{gumbel_type_of_selfSimilar_index_zero', gumbel_type_of_selfSimilar_index_zero}
  \leanok
  Suppose that $G$ is a nondegenerate c.d.f. such that
  \begin{align*}
    G^{t} = A_t . G \qquad \text{ for any } t > 0 ,
  \end{align*}
  where
  \begin{align*}
    A_t(x) = x + \beta \, \log t
  \end{align*}
  with $\beta > 0$.

  Then with $d = \log \big(-\log G(0) \big)$,
  for all $x \in \bR$ we have
  \begin{align*}
    G(x) = \exp \Big( - \exp \big( -\beta^{-1} x + d \big) \Big) .
  \end{align*}
  (In particular, $G$ is of Gumbel type: there exists an $A \in \OriAffR$
  such that $G = A.\GumbelCDF$.)
\end{lemma}
\begin{proof}
  %\uses{lem:solution-functional-eqn-translation-no-scaling}
  \uses{lem:no-fixed-point-implies-translation}
  % \leanok

  Since $G$ is nondegenerate, there exists an $x_0 \in \bR$
  with $0 < G(x_0) < 1$. Write $q = -\log G(x_0) > 0$,
  so that $G(x_0) = e^{-q}$.

  For $t > 0$, from the equation $G^{t} = A_t . G$
  we get for any $x \in \bR$
  \begin{align*}
  G(x)^t = G \big( A_t^{-1}(x) \big) = G \big( x - \beta \, \log(t) \big) .
  \end{align*}
  In particular, with $x = x_0 + \beta \, \log(t)$, we get
  \begin{align*}
    G\big( x_0 + \beta \, \log(t) \big)^{t} = G(x_0) = e^{- q},
  \end{align*}
  from which we can solve
  \begin{align*}
    G\big( x_0 + \beta \, \log(t) \big) = e^{-q/t} .
  \end{align*}
  The above holds for any $t>0$, and any $x \in \bR$
  can be written as $x = x_0 + \beta \, \log\big( e^{(x - x_0)/\beta} \big)$.
  We therefore get, for any $x \in \bR$,
  \begin{align*}
    G(x) = \exp \Big( - q \, \exp \big( -(x-x_0)/\beta \big) \Big) .
  \end{align*}
  This is of the desired form, with $d = \frac{x_0}{\beta} + \log(q)$.
  Plugging in $x = 0$ then shows $d = \log \big(-\log G(0) \big)$.
\end{proof}

\begin{lemma}[Self-similar continuous c.d.f. family characterization $\gamma > 0$]
  \label{lem:characterization-self-similar-family-pos-index}
  \uses{def:lem:oriented-affine-action-on-cdf, def:fixing-subgroup}
  \lean{frechet_type_of_selfSimilar_index_pos', frechet_type_of_selfSimilar_index_pos}
  \leanok
  Suppose that $G$ is a nondegenerate c.d.f. such that
  \begin{align*}
    G^{t} = A_t . G \qquad \text{ for any } t > 0 ,
  \end{align*}
  where
  \begin{align*}
    A_t(x) = t^{\alpha} x + c \, (1 - t^{\alpha}) = t^{\alpha} (x-c) + c
  \end{align*}
  with $c \in \bR$ and $\alpha > 0$.

  Then with $\sigma = \big(- \log G(c+1)\big)^{\alpha}$,
  for all $x \ge c$ we have
  \begin{align*}
    G(x) = \exp \Big( - \big( \frac{x - c}{\sigma} \big)^{-1/\alpha} \big) \Big) .
  \end{align*}

  (It easily follows that $G$ is of Fr\'echet type: there exists
  an $A \in \OriAffR$ such that $G = A.\FrechetCDF{\alpha}$.)
\end{lemma}
\begin{proof}
  %\uses{lem:solution-functional-eqn-translation-with-scaling}
  \uses{lem:fixing-subgroup-characterization, lem:cdf-equal-on-dense}
  % \leanok

  Since $G$ is nondegenerate, there exists an $x_0 \in \bR$
  with $0 < G(x_0) < 1$. Write $q = -\log G(x_0) > 0$,
  so that $G(x_0) = e^{-q}$.

  %Note that
  %\begin{align*}
  %  A_t^{-1}(x) = t^{-\alpha} x + c (1 - t^{-\alpha}) = t^{-\alpha} (x-c) + c .
  %\end{align*}
  For $t > 0$, from the equation $G^{t} = A_t . G$
  we get for any $x \in \bR$
  \begin{align*}
  G(x)^t = G \big( A_t^{-1}(x) \big) . %= G \big( t^{-\alpha} (x-c) + c \big) .
  \end{align*}

  Note that for $x \le c$ and $t = 2$ we have $A_2^{-1}(x) = 2^{-\alpha}(x-c) + c \ge x$,
  so the above implies $G(x)^2 = G(A_2^{-1}(x)) \ge G(x)$.
  This not possible if $0 < G(x) < 1$, so we must in particular have $x_0 > c$.

  With $x = A_t(x_0)$ in the above equation, %= t^{\alpha} (x_0-c) + c$,
  we get
  \begin{align*}
    G\big( A_t(x_0) \big)^{t} = G(x_0) = e^{- q},
  \end{align*}
  from which we can solve
  \begin{align*}
    G(x) = G\big( A_t(x_0) \big) = e^{-q/t} .
  \end{align*}

  %The above holds for any $t>0$.
  Any $x \ge c$ can be written as $x = A_t(x_0) = t^{\alpha} (x_0-c) + c$
  with $t = \big( \frac{x - c}{x_0 - c} \big)^{1/\alpha}$
  (recall that $x_0 - c > 0$).
  We therefore get, for any $x \ge c$,
  \begin{align*}
    G(x) = \exp \Big( - q \, \big( \frac{x - c}{x_0 - c} \big)^{-1/\alpha} \big) \Big) .
  \end{align*}
  This is of the form
  \begin{align*}
    G(x) = \exp \Big( - \big( \frac{x - c}{\sigma} \big)^{-1/\alpha} \big) \Big) ,
  \end{align*}
  and plugging in $x = c + 1$ yields $\sigma = \big(- \log G(c+1)\big)^{\alpha}$.
  %This is of the desired form, with $d = \frac{x_0}{c} + \log(q)$.
  %Plugging in $x = 0$ then shows $d = \log \big(-\log G(0) \big)$.
\end{proof}

\begin{lemma}[Self-similar continuous c.d.f. family characterization $\gamma < 0$]
  \label{lem:characterization-self-similar-family-neg-index}
  \uses{def:lem:oriented-affine-action-on-cdf, def:fixing-subgroup}
  \lean{weibull_type_of_selfSimilar_index_neg', weibull_type_of_selfSimilar_index_neg}
  \leanok
  Suppose that $G$ is a nondegenerate c.d.f. such that
  \begin{align*}
    G^{t} = A_t . G \qquad \text{ for any } t > 0 ,
  \end{align*}
  where
  \begin{align*}
    A_t(x) = t^{-\alpha} x + c \, (1 - t^{-\alpha}) = t^{-\alpha} (x-c) + c
  \end{align*}
  with $c \in \bR$ and $\alpha > 0$.

  Then with $\sigma = \big(- \log G(c-1)\big)^{-\alpha}$,
  for all $x \le c$ we have
  \begin{align*}
    G(x) = \exp \Big( - \big( \frac{c - x}{\sigma} \big)^{1/\alpha} \big) \Big) .
  \end{align*}

  (It easily follows that $G$ is of Weibull type: there exists
  an $A \in \OriAffR$ such that $G = A.\WeibullCDF{\alpha}$.)
\end{lemma}
\begin{proof}
  %\uses{lem:solution-functional-eqn-scaling,
  %  lem:solution-functional-eqn-translation-exp-zero}
  \uses{lem:fixing-subgroup-characterization, lem:cdf-equal-on-dense}
  % \leanok

  \emph{(Note: The Lean formalized statement uses the opposite sign of $\alpha$:
  it is assumed that $\alpha < 0$ and $A_t(x) = t^{+\alpha} (x-c) + c$.)}

  Since $G$ is nondegenerate, there exists an $x_0 \in \bR$
  with $0 < G(x_0) < 1$. Write $q = -\log G(x_0) > 0$,
  so that $G(x_0) = e^{-q}$.

  %Note that
  %\begin{align*}
  %  A_t^{-1}(x) = t^{-\alpha} x + c (1 - t^{-\alpha}) = t^{-\alpha} (x-c) + c .
  %\end{align*}
  For $t > 0$, from the equation $G^{t} = A_t . G$
  we get for any $x \in \bR$
  \begin{align*}
  G(x)^t = G \big( A_t^{-1}(x) \big) . %= G \big( t^{-\alpha} (x-c) + c \big) .
  \end{align*}

  Note that for $x \ge c$ and $t = 2$ we have $A_2^{-1}(x) = 2^{\alpha}(x-c) + c \ge x$,
  so the above implies $G(x)^2 = G(A_2^{-1}(x)) \ge G(x)$.
  This not possible if $0 < G(x) < 1$, so we must in particular have $x_0 < c$.

  With $x = A_t(x_0)$ in the above equation, %= t^{\alpha} (x_0-c) + c$,
  we get
  \begin{align*}
    G\big( A_t(x_0) \big)^{t} = G(x_0) = e^{- q},
  \end{align*}
  from which we can solve
  \begin{align*}
    G(x) = G\big( A_t(x_0) \big) = e^{-q/t} .
  \end{align*}

  %The above holds for any $t>0$.
  Any $x \le c$ can be written as $x = A_t(x_0) = t^{-\alpha} (x_0-c) + c$
  with $t = \big( \frac{c - x}{c - x_0} \big)^{-1/\alpha}$
  (recall that $c - x_0 > 0$).
  We therefore get, for any $x \le c$,
  \begin{align*}
    G(x) = \exp \Big( - q \, \big( \frac{c - x}{c - x_0} \big)^{1/\alpha} \big) \Big) .
  \end{align*}
  This is of the form
  \begin{align*}
    G(x) = \exp \Big( - \big( \frac{c - x}{\sigma} \big)^{1/\alpha} \big) \Big) ,
  \end{align*}
  and plugging in $x = c - 1$ yields $\sigma = \big(- \log G(c-1)\big)^{-\alpha}$.
  %This is of the desired form, with $d = \frac{x_0}{c} + \log(q)$.
  %Plugging in $x = 0$ then shows $d = \log \big(-\log G(0) \big)$.
\end{proof}


\begin{theorem}[Three types of extreme value distributions [Fisher-Tippett-Gnedenko]]
  \label{thm:three-types-of-extr-val-distr}
  \uses{def:extr-val-distr,
    def:std-Gumbel-cdf, def:std-Frechet-cdf, def:std-Weibull-cdf}
  \lean{CumulativeDistributionFunction.IsExtremeValueDistr.classification}
  \leanok
  For any extreme value distribution $G$, one of the following holds:
  \begin{description}
    \item{($\GumbelCDF$)}
      $G = A . \GumbelCDF$ for some $A \in \OriAffR$;
    \item{($\FrechetCDF{{}}$)}
      $G = A . \FrechetCDF{\alpha}$ for some $A \in \OriAffR$ and $\alpha > 0$;
    \item{($\WeibullCDF{{}}$)}
      $G = A . \WeibullCDF{\alpha}$ for some $A \in \OriAffR$ and $\alpha > 0$.
  \end{description}
  In particular, the only three possible types of extreme value distributions
  are the type of the Gumbel c.d.f., the type of
  the Fr\'echet c.d.f. $\FrechetCDF{\alpha}$ for $\alpha > 0$, and
  the type of the Weibull c.d.f. $\WeibullCDF{\alpha}$ for $\alpha > 0$.
\end{theorem}
\begin{proof}
  \uses{thm:one-parameter-subgroups-of-affine-isomorphisms,
    lem:self-similarity-of-extreme-value-distributions,
    lem:characterization-self-similar-family-zero-index,
    lem:characterization-self-similar-family-pos-index,
    lem:characterization-self-similar-family-neg-index}
  \leanok
  \ldots
\end{proof}

